\documentclass[10pt,a4paper]{article}
\usepackage{amsmath,amssymb,amsthm,graphicx}
\usepackage{amsmath,amssymb,amsthm,amsfonts}
\usepackage{latexsym,bbm,graphicx,float,mathtools}
\usepackage{subcaption}

\usepackage{verbatim}
\usepackage{enumitem}
\usepackage{hyperref}  
\hypersetup{
    colorlinks,
    citecolor=black,
    filecolor=black,
    linkcolor=black,
    urlcolor=black,
    linktoc=all
}

\usepackage[noend]{algpseudocode}
\usepackage{algorithm}

\usepackage{thmtools}
\usepackage{thm-restate}

%primed-thm
\newtheorem{manualtheoreminner}{Theorem}
\newenvironment{manualtheorem}[1]{%
  \renewcommand\themanualtheoreminner{#1}%
  \manualtheoreminner
}{\endmanualtheoreminner}

%footnote twice preamble
%https://tex.stackexchange.com/questions/35043/reference-different-places-to-the-same-footnote
\makeatletter
\newcommand\footnoteref[1]{\protected@xdef\@thefnmark{\ref{#1}}\@footnotemark}
\makeatother



 % environments
\newtheorem{problem}{Problem}
\newtheorem{theorem}{Theorem}[section]
\newtheorem{lem}[theorem]{Lemma}

\newtheorem{proposition}[theorem]{Proposition}
\newtheorem{fact}[theorem]{Fact}
\newtheorem{cor}[theorem]{Corollary}
\newtheorem{conjecture}[theorem]{Conjecture}
\newtheorem{question}{Question}
\newtheorem{goal}{Goal}

% non-italicized environments
\theoremstyle{definition}%\newtheorem{definition}[theorem]{Definition}
%\newtheorem{remark}[theorem]{Remark}
\newtheorem{remark}[theorem]{Remark}
\newtheorem{claim}[theorem]{Claim}
\newtheorem{definition}[theorem]{Definition}
\newtheorem{obs}[theorem]{Observation}
\newtheorem{notation}[theorem]{Notation}
\newtheorem{example}[theorem]{Example}
\newtheorem{examples}[theorem]{Examples}
\theoremstyle{plain}

% probability
\newcommand{\std}{\operatorname{{\bf std}}}
\newcommand{\Var}{\operatorname{{\bf Var}}}
\newcommand{\Cov}{\operatorname{{\bf Cov}}}
\newcommand{\Varx}{\mathop{{\bf Var}\/}}
%\newcommand{\E}{\operatorname{{\mathbb E}}}
\DeclareMathOperator*{\E}{\mathbb{E}}
%\renewcommand{\Pr}{\operatorname{{\mathbb P}}}
\let\Pr\relax
\DeclareMathOperator*{\Pr}{\mathbb{P}}
\newcommand{\Bin}{\operatorname{Bin}}
\newcommand{\Ber}{\operatorname{Ber}}

% mathrm terms
\newcommand{\sbp}{\sum\mathrm{smABP}}
\newcommand{\poly}{\mathrm{poly}}
\newcommand{\imm}[2]{\mathrm{IMM}_{#1,#2}}
\newcommand{\rk}{\mathrm{relrk}}
\newcommand{\polylog}{\mathrm{polylog}}
\newcommand{\size}{\mathrm{size}}
\newcommand{\avg}{\mathop{\mathrm{avg}}}
\newcommand{\odd}{\mathrm{odd}}
\newcommand{\even}{\mathrm{even}}
\newcommand{\sign}{\mathrm{sign}}
\newcommand{\dist}{\mathrm{dist}}
\newcommand{\dtv}{d_\mathrm{TV}}
\newcommand{\supp}{\mathrm{supp}}
\newcommand{\Vol}{\mathrm{Vol}}
\newcommand{\opt}{\mathrm{opt}}
\newcommand{\rank}{\mathrm{rank}}
\newcommand{\Tr}{\mathrm{Tr}}
\newcommand{\Alg}{\mathrm{Alg}}
\newcommand{\sen}{\mathrm{sen}}
\newcommand{\dnf}{\mathrm{dnf}}
\newcommand{\dt}{\mathrm{dt}}
\newcommand{\type}{\mathrm{type}}
\newcommand{\cent}{\mathrm{center}}
\newcommand{\dyck}{\mathrm{Dyck}}
\newcommand{\out}{\mathrm{out}}
\newcommand{\children}{\mathrm{children}}
\newcommand{\Cop}{\mathsf{Copies}}

%mathsf complexity classes
\newcommand{\VP}{\mathsf{VP}}
\newcommand{\VBP}{\mathsf{VBP}}
\newcommand{\VNP}{\mathsf{VNP}}
\newcommand{\VF}{\mathsf{VF}}
%\newcommand{\P}{\mathsf{P}}
\newcommand{\NP}{\mathsf{NP}}




% mathbf terms
\newcommand{\Size}{\mathbf{Size}}
\newcommand{\Depth}{\mathbf{Depth}}

% number systems
\newcommand{\R}{\mathbb R}
\newcommand{\C}{\mathbb C}
\newcommand{\N}{\mathbb N}
\newcommand{\Z}{\mathbb Z}
\newcommand{\F}{\mathbb F}
\newcommand{\K}{\mathbb K}
\newcommand{\Q}{\mathbb Q}

% short forms
\newcommand{\eps}{\varepsilon}
\renewcommand{\epsilon}{\eps} 
\newcommand{\la}{\langle}
\newcommand{\ra}{\rangle}
\newcommand{\wt}[1]{\widetilde{#1}}
\newcommand{\wh}[1]{\widehat{#1}}
\newcommand{\sse}{\subseteq}
\newcommand{\bits}{\{-1,1\}}
\newcommand{\zo}{\{0,1\}}
\newcommand{\bn}{\bits^n}
\newcommand{\isafunc}{:\bn\to\bits}
\newcommand{\isazofunc}{:\zo^n\to\zo}
\newcommand{\sumi}{\sum_{i=1}^n}

\newcommand{\Fsm}{\mathbb{F}_{\mathrm{sm}}}
\newcommand{\Ksm}{\mathbb{K}_{\mathrm{sm}}}
\newcommand{\lfrac}{\textstyle \frac}



% Overline
\newcommand{\ol}[1]{\overline{#1}}

% BOLD
\newcommand{\ind}{{\bf 1}}

\newcommand{\bA}{\mathbf{A}}
\newcommand{\bB}{\mathbf{B}}
\newcommand{\bC}{\mathbf{C}}
\newcommand{\bD}{\mathbf{D}}
\newcommand{\bE}{\mathbf{E}}
\newcommand{\bF}{\mathbf{F}}
\newcommand{\bG}{\mathbf{G}}
\newcommand{\bH}{\mathbf{H}}
\newcommand{\bI}{\mathbf{I}}
\newcommand{\bJ}{\mathbf{J}}
\newcommand{\bK}{\mathbf{K}}
\newcommand{\bL}{\mathbf{L}}
\newcommand{\bM}{\mathbf{M}}
\newcommand{\bN}{\mathbf{N}}
\newcommand{\bO}{\mathbf{O}}
\newcommand{\bP}{\mathbf{P}}
\newcommand{\bQ}{\mathbf{Q}}
\newcommand{\bR}{\mathbf{R}}
\newcommand{\bS}{\mathbf{S}}
\newcommand{\bT}{\mathbf{T}}
\newcommand{\bU}{\mathbf{U}}
\newcommand{\bV}{\mathbf{V}}
\newcommand{\bW}{\mathbf{W}}
\newcommand{\bX}{\mathbf{X}}
\newcommand{\bY}{\mathbf{Y}}
\newcommand{\bZ}{\mathbf{Z}}


\newcommand{\boldf}{\boldsymbol{f}} % \bf taken
\newcommand{\ba}{\boldsymbol{a}}
\newcommand{\bb}{\boldsymbol{b}}
\newcommand{\be}{\boldsymbol{e}}
\newcommand{\bg}{\boldsymbol{g}}
\newcommand{\bh}{\boldsymbol{h}}
\newcommand{\bj}{\boldsymbol{j}}
\newcommand{\bk}{\boldsymbol{k}}
\newcommand{\bm}{\boldsymbol{m}}
\newcommand{\bp}{\boldsymbol{p}}
\newcommand{\boldr}{\boldsymbol{r}} % \br taken
\newcommand{\bs}{\boldsymbol{s}}
\newcommand{\bt}{\boldsymbol{t}}
\newcommand{\bu}{\boldsymbol{u}}
\newcommand{\bv}{\boldsymbol{v}}
\newcommand{\bw}{\boldsymbol{w}}
\newcommand{\bx}{\boldsymbol{x}}
\newcommand{\by}{\boldsymbol{y}}
\newcommand{\bz}{\boldsymbol{z}}

\newcommand{\boldR}{\boldsymbol{R}}


\newcommand{\bsigma}{\boldsymbol{\sigma}}
\newcommand{\bpi}{\boldsymbol{\pi}}
\newcommand{\bell}{\boldsymbol{\ell}}
\newcommand{\bnu}{\boldsymbol{\nu}}
\newcommand{\bxi}{\boldsymbol{\xi}}
\newcommand{\btau}{\boldsymbol{\tau}} 
\newcommand{\brho}{{\boldsymbol{\rho}} }
\newcommand{\bgamma}{{\boldsymbol{\gamma}} }
\newcommand{\bphi}{{\boldsymbol{\varphi}} }


\newcommand{\bvsigma}{{\boldsymbol{\varsigma}} }
\newcommand{\bvrho}{{\boldsymbol{\varrho}}}

\newcommand{\bGamma}{{\mathbf{\Gamma}} }


% Calligraphy

\newcommand{\calA}{\mathcal{A}}
\newcommand{\calB}{\mathcal{B}}
\newcommand{\calC}{\mathcal{C}}
\newcommand{\calD}{\mathcal{D}}
\newcommand{\calE}{\mathcal{E}}
\newcommand{\calF}{\mathcal{F}}
\newcommand{\calG}{\mathcal{G}}
\newcommand{\calH}{\mathcal{H}}
\newcommand{\calI}{\mathcal{I}}
\newcommand{\calJ}{\mathcal{J}}
\newcommand{\calK}{\mathcal{K}}
\newcommand{\calL}{\mathcal{L}}
\newcommand{\calM}{\mathcal{M}}
\newcommand{\calN}{\mathcal{N}}
\newcommand{\calO}{\mathcal{O}}
\newcommand{\calP}{\mathcal{P}}
\newcommand{\calQ}{\mathcal{Q}}
\newcommand{\calR}{\mathcal{R}}
\newcommand{\calS}{\mathcal{S}}
\newcommand{\calT}{\mathcal{T}}
\newcommand{\calU}{\mathcal{U}}
\newcommand{\calV}{\mathcal{V}}
\newcommand{\calW}{\mathcal{W}}
\newcommand{\calX}{\mathcal{X}}
\newcommand{\calY}{\mathcal{Y}}
\newcommand{\calZ}{\mathcal{Z}}

\newenvironment{solution}
  {\begin{proof}[\textbf{Solution}]}
  {\end{proof}}
 
\usepackage{cleveref}
\crefname{originalsection}{original section}{original sections}
\Crefname{originalsection}{Original section}{Original sections}
\makeatletter
\newlabel{subsec:our-results}{{1.3}{source}{}{sourcepointer-1}{}}
\newlabel{subsec:our-results@cref}{{[originalsection][0][]1.3}{[1][0][]source}{}{}{}}
\makeatother
\crefname{originaltheorem}{original theorem}{original theorems}
\Crefname{originaltheorem}{Original theorem}{Original theorems}
\makeatletter
\newlabel{thm:large-deg-VP}{{1.6}{source}{}{sourcepointer-2}{}}
\newlabel{thm:large-deg-VP@cref}{{[originaltheorem][0][]1.6}{[1][0][]source}{}{}{}}
\makeatother
\crefname{originaltheorem}{original theorem}{original theorems}
\Crefname{originaltheorem}{Original theorem}{Original theorems}
\makeatletter
\newlabel{thm:small-deg-VP}{{1.4}{source}{}{sourcepointer-3}{}}
\newlabel{thm:small-deg-VP@cref}{{[originaltheorem][0][]1.4}{[1][0][]source}{}{}{}}
\makeatother
\crefname{originaltheorem}{original theorem}{original theorems}
\Crefname{originaltheorem}{Original theorem}{Original theorems}
\makeatletter
\newlabel{thm:tightness}{{1.8}{source}{}{sourcepointer-4}{}}
\newlabel{thm:tightness@cref}{{[originaltheorem][0][]1.8}{[1][0][]source}{}{}{}}
\makeatother
\title{Lower bounds for sums of ordered set-multilinear ABPs computing a polynomial with a small general ABP}
\author{Prerona Chatterjee, Deepanshu Kush, Shubhangi Saraf, Amir Shpilka}
\date{doi:10.4230/LIPIcs.CCC.2024.20; CC BY-NC-SA4.0}
\begin{document}
\maketitle
\noindent\textit{Editorial note.} Original source: Prerona Chatterjee, Deepanshu Kush, Shubhangi Saraf and Amir Shpilka; arXiv:2312.15874v3 (2024-02-19), https://arxiv.org/abs/2312.15874v3. License: Creative Commons Attribution-NonCommercial-ShareAlike4.0, https://creativecommons.org/licenses/by-nc-sa/4.0/.\par\medskip
\noindent\textit{Editorial note.} Non-commercial research collection. Changes: complete author theorem/proof/context excerpts and original macros/bibliography in standalone layout; unrelated sections and frontmatter omitted. No mathematical corrections or additional proofs. Adapted presentation retains CC BY-NC-SA4.0 including NC/SA restrictions.\par\medskip
\noindent\textit{Allowed standard prerequisites (editorial).}\begin{itemize}
\item Standard relative-rank properties and inner-product gadget. Author Claim clm:rk-props and cited Observation obs:ip-full-rank; precise statements retained.
\item Standard Chernoff bound. Author Theorem thm:Chernoff explicitly cites Mitzenmacher-Upfal; exact statement retained.
\end{itemize}\medskip
\ifdefined\hypertarget\hypertarget{sourcepointer-1}{}\fi
\noindent\textit{Source reference.} Original section 1.3: number-theory/sources/ccc-2024-20/native/main.tex, source lines 150--150. Optional introductory pointer to results; counted full target statements and proofs are included.\par\medskip
\ifdefined\hypertarget\hypertarget{sourcepointer-2}{}\fi
\noindent\textit{Source reference.} Original theorem 1.6: number-theory/sources/ccc-2024-20/native/main.tex, source lines 176--176. Only supplies the source numbering for the primed theorem; the complete VBP theorem/proof is included.\par\medskip
\ifdefined\hypertarget\hypertarget{sourcepointer-3}{}\fi
\noindent\textit{Source reference.} Original theorem 1.4: number-theory/sources/ccc-2024-20/native/main.tex, source lines 154--154. Only supplies the source numbering for the primed theorem; complete counted VBP target is included.\par\medskip
\ifdefined\hypertarget\hypertarget{sourcepointer-4}{}\fi
\noindent\textit{Source reference.} Original theorem 1.8: number-theory/sources/ccc-2024-20/native/main.tex, source lines 203--203. Author mentions this as an immediate consequence of the delivered construction and proof; it is not a premise or separate counted item.\par\medskip
An \emph{algebraic circuit} over a field $\F$ for a multivariate polynomial $P(x_1,\ldots,x_N)$ is a directed acyclic graph (DAG)
whose internal vertices (called gates) are labeled as either $+$ (sum) or $\times$ (product), and leaves
(vertices of in-degree zero) are labeled by the variables $x_i$ or constants from $\F$. 
A special output gate (the root of the DAG) represents the polynomial
$P$. 
If the DAG happens to be a tree, such a resulting circuit is called an \emph{algebraic formula}. 
The size of a circuit or formula is the number of nodes in the DAG. We also consider the product-depth of the circuit, which is the maximum number of product gates on a root-to-leaf path. 
The class $\VP$ (respectively, $\VF$) is then defined to be the collection of all polynomials having at most polynomially large degree which can be computed by polynomial-sized circuits (respectively, formulas).

An \emph{algebraic branching program} (ABP) is a layered DAG with two special nodes in
it: a start-node and an end-node. All edges of the ABP go from layer $\ell -1$ to layer $\ell$ for some $\ell$ (say start-node is the unique node in layer $0$ and end-node is the unique node in the last layer) and are labeled by a linear polynomial. Every directed path $\gamma$ from start-node to end-node computes the monomial $P_\gamma$,
which is the product of all labels on the path $\gamma$. The ABP computes the polynomial $P = \sum_{\gamma} P_\gamma$,
where the sum is over all paths $\gamma$ from start-node to end-node. Its size is simply the number of nodes in the DAG, its \emph{depth} is the length of the longest path from the start-node to the end-node, and \emph{width} is the maximum number of nodes in any layer. The class $\VBP$ is then defined to be the collection of all polynomials (with polynomially-bounded degree) which can be computed by polynomial-sized branching programs. ABPs are known to be of intermediate complexity between formulas and circuits; in other words, we know the inclusions $\VF\subseteq\VBP \subseteq \VP\sse \VNP$.

A polynomial is said to be \emph{homogeneous}
if each monomial has the same total degree and \emph{multilinear} if every variable occurs
at most once in any monomial. Now, suppose that the underlying variable set is partitioned into $d$ sets
$X_1,\ldots, X_d$. Then the polynomial is said to be \emph{set-multilinear} with respect to this variable partition if each
monomial in $P$ has \emph{exactly} one variable from each set. Note that a set-multilinear polynomial is both multilinear and homogeneous, and has degree precisely $d$ if it is set-multilinear over $d$ sets. Next, we define different models
of computation corresponding to these variants of polynomials classes. An algebraic formula/branching program/circuit is set-multilinear with respect to a variable partition
$(X_1,\ldots, X_d)$ if each internal node in the formula/branching program/circuit computes a set-multilinear polynomial.\footnote{Of course, a non-root node need not be set-multilinear with respect to the \emph{entire} variable partition. Nevertheless, here we demand that it must be set-multilinear with respect to some \emph{subset} of the collection $\{X_1,\ldots,X_d\}$.} Multilinear and homogeneous formulas/branching programs/circuits are defined analogously.  

\begin{definition}[Ordered smABP]\label{def:smABP}
Given a variable partition $(X_1,\ldots, X_d)$, we say that a set-multilinear branching program of depth $d$ is said to be \emph{ordered} with respect to an 
ordering (or permutation) $\sigma\in S_d$ if for each $\ell\in[d]$, all edges of the ABP from layer $\ell-1$ to layer $\ell$ are labeled using a linear form over the variables in $X_{\sigma(\ell)}$. It is simply said to be \emph{ordered} if there \emph{exists} an ordering $\sigma$ such that it is ordered with respect to $\sigma$.
\end{definition}

At this point, it is essential to take note of the terminology in this context: in this paper, a general (or ``unordered'') set-multilinear branching program refers to an ABP for which each internal node computes a polynomial that is set-multilinear with respect to \emph{some} subset of the global partition, whereas an \emph{ordered} set-multilinear branching program is more specialized and has the property that any two nodes in the same layer compute polynomials that are set-multilinear with respect to the \emph{same} partition.

\begin{definition}[$\sbp$]\label{def:sum}
Given a polynomial $P(X)$ that is set-multilinear with respect to the variable partition $X = (X_1,\ldots, X_d)$, we say that $\sum_{i=1}^t A_i$ is a $\sbp$ computing $P$ if indeed $\sum_{i=1}^t A_i(X) = P(X)$, and
each $A_i$ is an ordered set-multilinear branching program i.e., each $A_i$ is ordered with respect to some $\sigma_i\in S_d$. We call $t$ (i.e., the number of summands in a $\sbp$) its \emph{support size} and define its \emph{max-width} and \emph{total-width} to be the maximum over the width of each $A_i$ and the sum of the width of each $A_i$, respectively.    
\end{definition}

\begin{manualtheorem}{\ref{thm:large-deg-VP}'}\label{thm:large-deg-VBP}
    There is a fixed constant $\delta \geq 1/100$ and a set-multilinear polynomial $G_{n,n}$ in $\VBP$, in  $\Theta(n^2)$ variables and of degree $\Theta(n)$,  such that any $\sbp$ computing $G_{n,n}$ requires total-width ${\exp(\Omega(n^{\delta}))}$.
\end{manualtheorem}

\begin{manualtheorem}{\ref{thm:small-deg-VP}'}\label{thm:small-deg-VBP}
Let $d\leq n$ be growing parameters satisfying $d = \omega(\log n)$. There is a  $\Theta(dn)$-variate, degree $\Theta(d)$ set-multilinear polynomial $G_{n,d}$ in $\VBP$ such that $G_{n,d}$  cannot be computed by a $\sbp$ of total-width $\poly(n)$.
\end{manualtheorem}

\section{Preliminaries}\label{sec:prel}
\subsection{Relative Rank and its Properties}

We first describe the notation that we need to define the measures that we use to prove our results described in Section \ref{subsec:our-results}. Instead of directly working with the rank of the partial derivative matrix, we work with the following normalized form.

\begin{definition}\label{def:rk-rest}
Let $w = (w_1, w_2,\ldots, w_d)$ be a tuple (or word) of non-zero real numbers. For a subset $S \subseteq [t]$, we shall refer to the sum $\sum_{i\in S} w_i$ by $w_S$,  and by $w|_S$, we will refer to the tuple obtained by considering only the
elements of $w$ that are indexed by $S$. 
Given a word $w=(w_1,\ldots,w_d)$, we denote by $\ol{X}(w)$ a tuple of $d$ sets of variables $(X(w_1),\ldots,X(w_d))$ where
$|X(w_i)| = 2^{|w_i|}$.\footnote{In particular, $2^{|w_i|}\in\mathbb{N}$.}
We denote by $\Fsm[\calT]$ the set of set-multilinear polynomials over the tuple
of sets of variables $\calT$.
\end{definition}


\begin{definition}[Relative Rank Measure of \cite{LST1}]\label{def:relrk}
Let $\ol{X}=(X_1,\ldots,X_d)$ be a tuple of sets of variables such that $|X_i|=n_i$ and let $f\in \Fsm[\ol{X}]$.
 Let $w = (w_1, w_2,\ldots, w_d)$ be a tuple (or word) of non-zero real numbers such that $2^{|w_i|}  = n_i $ for all $i \in [d]$. Corresponding to a word $w$, define $\calP_w \coloneqq \{i\ |\ w_i > 0\}$ and $\calN_w \coloneqq \{i\ |\ w_i < 0\}$. Let $\calM^{\calP}_{w}$ be the
set of all set-multilinear monomials over the subset of the variable sets $X_1, X_2,\ldots, X_d$ precisely
indexed by $\calP_w$, and similarly let $\calM^{\calN}_{w}$ be the set of all set-multilinear monomials over these variable
sets indexed by $\calN_w$.

Define the ‘partial derivative matrix’ matrix $\calM_w(f)$ whose rows are indexed by the elements
of $\calM^{\calP}_w$ and columns indexed by the elements of $\calM^{\calN}_{w}$ as follows: the entry of this matrix corresponding to a
row $m_1$ and a column $m_2$ is the coefficient of the monomial $m_1\cdot m_2$ in $f$. We define
\[
\rk_w(f) \coloneqq \frac{\mathrm{rank}(\calM_w(f))}{\sqrt{|\calM^{\calP}_w|\cdot |\calM^{\calN}_w|}} = \frac{\mathrm{rank}(\calM_w(f))}{2^{\frac{1}{2}\sum_{i\in[d]}|w_i|}}.
\]
\end{definition}


The following is a simple result that establishes various useful properties of the relative rank measure.
\begin{claim}[\cite{LST1}]\label{clm:rk-props}
\begin{enumerate}
    \item(Imbalance) Say $f \in \Fsm[\ol{X}(w)]$. Then, $\rk_w(f)\leq 2^{-|w_{[d]}|/2}$.
    \item(Sub-additivity) If $f,g \in \Fsm[\ol{X}(w)]$, then $\rk_w(f+g)\leq \rk_w(f)+\rk_w(g)$.
    \item(Multiplicativity) Say $f = f_1 f_2\cdots f_t $ and assume that for each $i\in [t]$, $f_i\in \Fsm[\ol{X}(w|_{S_i})]$, where $(S_1, \ldots, S_t)$ is a partition of $[d]$. Then
    \[
    \rk_w(f) = \prod_{i\in[t]} \rk_{w|_{S_i}}(f_i).
    \]
\end{enumerate}
\end{claim}

We will repeatedly make use of the following.

\begin{theorem}[Chernoff Bound, as stated in \cite{mitzenmacher_upfal_2005}]\label{thm:Chernoff}
Suppose $X_1, \ldots, X_n$ are independent random variables taking values in $\zo$. Let $X$ denote their sum and let $\mu = \E[X]$ denote the expected value of the sum. Then for any $\delta > 0$,
\begin{eqnarray*} \Pr[X\ge (1+\delta)\mu] &\le& \exp\left(-\frac{\delta^2\mu}{2+\delta}\right), \text{ if }0\leq \delta, \\
  \Pr[X\le (1-\delta)\mu] &\le& \exp\left(-\frac{\delta^2\mu}{2}\right), \text{ if }0< \delta <1, \\ 
  \Pr[|X-\mu|\geq \delta \mu]&\leq& 2\cdot \exp\left(-\frac{\delta^2\mu}{3}\right),\text{ if }0< \delta <1.\\
\end{eqnarray*}
\end{theorem}

\subsection{Inner Product Gadget}
The following observation is used crucially to construct the hard polynomials in $\VP$ as well as $\VBP$. 
\begin{obs}[\cite{KushS23}]\label{obs:ip-full-rank}
    Let $n = 2^k$ and $X_1 = \{x_{1,1},\ldots,x_{1,n}\}$ and $X_2 = \{x_{2,1},\ldots,x_{2,n}\}$ be two disjoint sets of variables. Then, for any \emph{symmetric} word $w\in\{k,-k\}^2$ (i.e., where $w_1 + w_2 = 0$) and for the inner product `gadget' $ f = X_1\cdot X_2 = \sumi x_{1,i}x_{2,i}$, $\rk_w(f) = 1$ i.e., $\calM_w(f)$ is full-rank.
\end{obs}


\subsection{A Hard Set-multilinear Polynomial in $\VP$}\label{sec:VP}
Let $d$ be an even integer and let $X = (X_1,\ldots, X_{d})$ be a collection of sets of variables where each $|X_i| = n$, and similarly, let $Y = (Y_1,\ldots, Y_{d})$ be a distinct collection of sets of variables where each $|Y_i| = n$. We shall refer to the $Y$-variables as the \emph{auxiliary} variables. For $i$ and $j \in \{1,\ldots,d\}$, let $X_i \cdot X_j$ denote the inner-product quadratic form $\sum_{k = 1}^n x_{ik}x_{jk}$. Here, we shall assume that $X_i = \{x_{i,1},\ldots, x_{i,n}\}$ and $Y_i = \{y_{i,1},\ldots, y_{i,n}\}$.

For two integers $i \in \N$ and $j \in \N$, we denote $[i, j] = \{k \in \N : i \leq k \text{ and } k \leq j\}$ and call such a set an \emph{interval}. For every interval $[i, j] \subseteq [d]$, we define a polynomial $f_{i,j}(X,Y)\in \Fsm[X_i,\ldots,X_j,Y_i,\ldots, Y_j]$ as follows:

\[
f_{i,j} = \begin{cases}
y_{i,j}y_{j,i}(X_i\cdot X_j) &\text{if $j = i+1$}\\
0 &\text{if $j-i$ is even}\\
y_{i,j}y_{j,i}(X_i\cdot X_j)\cdot f_{i+1,j-1} + \sum_{r = i+1}^{j-1} f_{i,r}f_{r+1,j} &\text{otherwise}\\
\end{cases}
\]

These $f_{i,j}$ in present form were defined in \cite{KushS23}, but were in turn inspired from an earlier work of Raz and Yehudayof (\cite{RazY08}) in the multilinear context. \cite{KushS23} shows that they have the following full-rank property that will be instrumental for us.
\begin{lem}[\cite{KushS23}]\label{lem:Fn-VP-props}
    Let $n = 2^k$ and $d\leq n$ be an even integer. Over any field $\F$ of characteristic zero, the polynomial $F_{n,d} = f_{1,d}\in\Fsm[X,Y]$ as defined above satisfies the following: For any $w\in\{-k,k\}^d$ with $w_{[d]} = 0$, $\calM_w(F_{n,d})$ is full-rank when viewed as a matrix over the field $\F(Y)$, the field of rational functions over the $Y$ variables.
\end{lem}


\subsection{A Hard Set-Multilinear Polynomial in $\VBP$}\label{subsec:VBP}
\subsubsection{Arc-partition Measure Description}
This subsection is adapted from Section 2 of \cite{DvirMPY12}. Let $n = 2^k$, $d\leq n$ be an even integer, and let $X = (X_1, X_2, \ldots, X_{d})$ be a collection of disjoint sets of $n$ variables each.
An \emph{arc-partition} will be a special kind of \emph{symmetric} word $w\in \{-k,k\}^d$ (i.e., a one-to-one map $\Pi$ from $X$ to $\{-k,k\}^d$). For the purpose of this subsection, the reader can even choose to think of the alphabet of $w$ as $\{-1,1\}$ (i.e., one `positive' and one `negative' value) -- we use $k,-k$ only to remain consistent with \autoref{def:relrk}.

Identify $X$ with the set $\{1, 2,\ldots, d\}$ in the natural way. Consider the $d$-cycle
graph, i.e., the graph with nodes $\{1, 2,\ldots, d\}$ and edges between $i$ and $i + 1$ modulo $d$. For
two nodes $i \neq j$ in the $d$-cycle, denote by $[i, j]$ the arc between $i, j$, that is, the set of nodes on the
path $\{i, i + 1, \ldots , j - 1, j\}$ from $i$ to $j$ in $d$-cycle.
First, define a distribution $\calD_P$ on a family of pairings (a list of disjoint pairs of nodes in the
cycle) as follows. A random pairing is constructed in $d/2$ steps. At the end of step $t \in [d/2]$, we
shall have a pairing $(P_1,\ldots , P_t)$ of the arc $[L_t, R_t]$. The size of $[L_t
, R_t]$ is always $2t$. The first pairing
contains only $P_1 = \{L_1, R_1\}$ with $L_1 = 1$ and $R_1 = 2$. Given $(P_1, \ldots, P_t)$ and $[L_t, R_t]$, define the
random pair $P_t+1$ (independently of previous choices) by 

$$
P_{t+1} = \begin{cases}
    \{L_t - 2, L_t - 1\} &\text{ with probability $1/3$}\\
    \{L_t - 1, R_t + 1\} &\text{ with probability $1/3$}\\
    \{R_t + 1, R_t + 2\} &\text{ with probability $1/3$}\\
\end{cases}
$$

Define
$$
[L_{t+1},R_{t+1}] = [L_t,R_t] \cup P_{t+1}.
$$

So, $L_{t+1}$ is either $L_t - 2$, $L_t - 1$ or $L_t$, each value is obtained with probability $1/3$, and similarly
(but not independently) for $R_{t+1}$.

The final pairing is
$P = (P_1, P_2, \ldots , P_{d/2}
)$.
Denote by $P \sim \calD_P$ a pairing distributed according to $\calD_P$.

Once a pairing $P$ has been obtained, a word $w\in \{-k,k\}^d$ is obtained by simply randomly assigning $+k$ and $-k$ to the indices of any pair $P_i$. More formally, for every
$t \in [d/2]$, if $P_t = \{i_t
, j_t\}$, let with probability $1/2$, independently of all other choices,
$$
w_{i_t} = +k \text{ and } w_{j_t} = -k,
$$
and with probability 1/2,
$$
w_{i_t} = -k \text{ and } w_{j_t} = +k.
$$
Denote by $w\sim \calD$ a word in $\{-1,1\}^n$ that is sampled using this procedure. We call such a word an \emph{arc-partition}. For a pair $P_t = \{i_t,j_t\}$, we refer to $i_t$ and $j_t$ as \emph{partners}. 

\begin{definition}[Arc-full-rank]\label{def:arc-full-rk}
We say that a polynomial $f$ that is set-multilinear over $X = (X_1, \ldots, X_{d})$ is \textbf{arc-full-rank} if for every arc-partition $w\in \{-k,k\}^d$, $\rk_w(f) = 1$.
\end{definition}


\subsubsection{Construction of an Arc-full-rank Polynomial}\label{subsubsec:arc-full-rank-poly}

Below, we describe a simple construction of a polynomial sized ABP that computes an arc-full-rank set-multilinear polynomial.
The high-level idea is to construct an ABP in which every path between start-node and end-node
corresponds to a specific execution of the random process which samples arc-partitions. Each node in the ABP corresponds to an arc $[L, R]$, which sends an edge to each of the nodes $[L - 2, R], [L -
1, R + 1]$ and $[L, R + 2]$. The edges have specially chosen labels that help guarantee full rank with respect to
every arc-partition. For simplicity of presentation, we allow the edges of the program to be labeled
by degree four set-multilinear polynomial polynomials over the corresponding subset of the variable partition. This assumption can be easily removed by replacing
each edge with a polynomial-sized ABP computing the corresponding degree four polynomial.

Formally, the nodes of the program are even-size arcs in the $d$-cycle, $d$ an even integer. The
start-node of the program is the empty arc $\emptyset$ and the end-node is the whole cycle $[d]$ (both are
“special” arcs). Let $X = (X_1,\ldots, X_{d})$ be a collection of sets of variables where each $|X_i| = n$, and similarly, let $Y = (Y_1,\ldots, Y_{d})$ be a distinct collection of sets of variables where each $|Y_i| = n$ (we shall refer to the $Y$-variables as \emph{auxiliary} variables). For $i$ and $j$ in $\{1,\ldots,d\}$, let $X_i \cdot X_j$ denote the inner-product quadratic form $\sum_{k = 1}^n x_{ik}x_{jk}$. Here, we shall assume that $X_i = \{x_{i,1},\ldots, x_{i,n}\}$ and $Y_i = \{y_{i,1},\ldots, y_{i,n}\}$.

Construct the branching program by connecting a node/arc of size $2t$ to three nodes/arcs of
size $2t + 2$. For $t = 1$, there is just one node $[1, 2]$, and the edge from start-node to it is labeled
$y_{1,2}y_{2,1}(X_0\cdot X_1)$. For $t > 1$, the node $[L, R] \supset [1, 2]$ of size $2t < d$ is connected to the three nodes:
$[L-2, R] , [L-1, R+ 1]$, and $[L, R+ 2]$. (It may be the case that the three nodes are the end-node.)
The edge labeling is: 
\begin{itemize}
    \item The edge  between $[L, R]$ and $[L-2, R]$ is labeled $y_{L-2,L-1}y_{L-1,L-2}(X_{L-2}\cdot X_{L-1})$.
    \item The
edge between $[L, R]$ and $[L - 1, R + 1]$ is labeled $y_{L-1,R+1}y_{R+1,L-1}(X_{L-1}\cdot X_{R+1})$.
    \item The edge between $[L, R]$
and $[L, R + 2]$ is labeled $y_{R+1,R+2}y_{R+2,R+1}(X_{R+1}\cdot X_{R+2})$.
\end{itemize}
  
Consider the ABP thus described, and the polynomial $G_{n,d}$ it computes. For every path $\gamma$
from start-node to end-node in the ABP, the list of edges along $\gamma$ yields a pairing $P$; every edge $e$
in $\gamma$ corresponds to a pair $P_e = \{i_e, j_e\}$ of nodes in $d$-cycle. Thus,

\begin{equation}\label{eq:poly-def}
  G_{n,d} = \sum_{\gamma} \prod_{e = \{i_e, j_e\} \in \gamma}y_{i_e,j_e}y_{j_e,i_e} \cdot (X_{i_e}\cdot X_{j_e}). 
\end{equation}

where the sum is over all paths $\gamma$ from start-node to end-node. 

\begin{remark}\label{rem:parings-paths}
There is in fact a one-to-one
correspondence between pairings $P$ and such paths $\gamma$ (this follows by induction on $t$). {Note that this is true only because pairings are tuples i.e., they are \emph{ordered} by definition. Otherwise, it is of course still possible to obtain the same \emph{set} of pairs in a given pairing using multiple different orderings.} The sum
defining $G_{n,d}$ can be thought of, therefore, as over pairings $P$.  
\end{remark}

The following statement summarizes the main useful property of $G_{n,d}$.

\begin{lem}[\cite{KushS23}]\label{lem:Fn-VBP-props}
Over any field $\F$ of characteristic zero, the polynomial $ G_{n,d}$ defined above is arc-full-rank as a set-multilinear polynomial in the variables $X$ over the field $\F(Y)$ of rational functions
in $Y$.
\end{lem}

\begin{proof}
    Let $w\sim \calD$ be an arc-partition. 
    We want to show that $\calM_w(G_{n,d})$ has full rank. 
    The arc-partition $w$ is defined from a pairing $P = (P_1,\ldots,P_{d/2})$ (though as discussed in \autoref{rem:parings-paths}, there could be multiple such $P$). 
    The pairing $P$ corresponds to a path $\gamma$ from start-node to end-node. 
    Consider the polynomial $f$ that is obtained by setting every $y_{i,j} = y_{j,i} = 0$ in $F$ such that $\{i,j\}$ is not a pair in $P$, and setting every $y_{i,j} = y_{j,i} = 1$ for every pair $\{i,j\}$ in $P$. 
    Then, it is easy to see that the only terms that survive in \autoref{eq:poly-def} correspond to paths (and in turn, pairings) which have the same underlying \emph{set} of pairs as $P$.
    As a consequence, $f$ is simply some non-zero constant times a polynomial which is full-rank (recall \autoref{obs:ip-full-rank}).
    $M_w(f)$ being full rank then implies that $M_w(G_{n,d})$ is also full-rank.
\end{proof}

\section{Separation between $\VP$ and $\sbp$}

In the theorem below, $F_{n,d}$ refers to the polynomial defined in Section \ref{sec:VP}. 
In this section, we prove \autoref{thm:large-deg-VP} and \autoref{thm:small-deg-VP} by first proving the following statement. 

\begin{lem}\label{lem:VP}
    Given large enough integers $d\leq n$, any $\sbp$ of max-width $s$ and support size $t$ computing $F_{n,d}(X_1,\ldots,X_d,Y_1\ldots,Y_d)$ must satisfy at least one of the following:
    \begin{itemize}
        \item either $t>e^{d/96}$,
        \item or $t\cdot s^q \geq n^{\frac{\sqrt{dq}}{20}}$, for any integer $q$ in the range $ [12\ln(2t\sqrt{d}), d/4]$.
    \end{itemize}  
\end{lem}

\begin{proof}
    Suppose that $t\leq e^{d/96}$ (so that the range in the theorem statement is indeed well-defined). 
    
    First, we observe that for any $\sbp$ computing $F_{n,d}(X,Y)$, we can view each summand as an ordered set-multilinear branching program with respect to only the $X$ variables. 
    In other words, by appropriately collapsing the layers labelled using the $Y$ variables, each summand is a set-multilinear branching program over the field $\F(Y)$ \emph{ordered} with respect to $(X_{\sigma(1)},\ldots, X_{\sigma(d)})$ for some permutation $\sigma\in S_d$. 
    It is easy to see that this collapsing process does not increase the width or the size of any summand in the branching program. 
    The edge labels do get altered however: the coefficients of the $X$ variables in any edge label can now be polynomials in the $Y$ variables (and therefore, field constants in $\F(Y)$).

    Let $A$ be such a set-multilinear branching program of width $s$ and depth $d$ that is \emph{ordered} with respect to $(X_1,\ldots, X_d)$\footnote{In general, $A$ may be ordered with respect to an arbitrary permutation $\sigma$, but the assumption that $\sigma$ is the identity permutation in the discussion that follows is without loss of generality.}. 
    Recall that this means that for each $\ell\in[d]$, all edges of the ABP from layer $\ell-1$ to layer $\ell$ are labeled using a linear form {in} $\mathit{X_\ell}$. 
Given a node $u$ in layer $i$ and a node $v$ in layer $j>i$ of $A$, define $g_{u,v}$ to be the polynomial computed by the ABP restricted to the layers $i+1,\ldots,j-1$ with the source and the sink defined by $u$ and $v$ respectively. Consider an integer $q$ in the range specified in the theorem statement and let $r$ be the largest  integer $r$ such that the product $qr< d$ i.e., we must have $d-qr\leq q$. Consider the following decomposition of $A$:
\[
A(X_1,\ldots,X_d) = \sum_{u_1,\ldots, u_{q}} \prod_{i = 1}^{q+1} g_{u_{i-1},u_i},
\]
where $u_0$ and $u_{q+1}$ are defined to be the source and the sink of $A$ respectively, and for $1\leq i \leq q$, $u_i$ varies over all choices of nodes in layer $r\cdot i$. Note that hence, this expression contains at most $s^q$ terms. Also, note that each $g_{u_{i-1},u_i}$ is set-multilinear over the partition $(X_i)_{i\in S_i}$ where $S_i$ is the set $\{r(i-1) + 1, \ldots, ri\}$\footnote{If $A$ is instead ordered with respect to $\sigma$, then $S_i$ is taken to be the set $\{r(i-1) + 1, \ldots, ri\}$.} of length exactly $r$ for $i\in [q]$, and $S_{q+1}=\{rq+1,\ldots,d\}$ has length at most $q$. We now analyze the relative rank of each summand. 

Let  $w\in \{-k,k\}^d$ (where $k = \log n$) be an arbitrary word. By \autoref{clm:rk-props}, we see that 

\[
\rk_w\left(\prod_{i = 1}^{q+1} g_{u_{i-1},u_i}\right) = \prod_{i = 1}^{q+1} \rk_{w|_{S_i}}(g_{u_{i-1},u_i})
\leq \prod_{i = 1}^{q} \rk_{w|_{S_i}}(g_{u_{i-1},u_i})\leq \prod_{i = 1}^q 2^{-|w_{S_i}|/2}= 2^{-\frac12\left(\sum_{i=1}^q |w_{S_i}|\right)},
\]
from which we observe that the task of upper bounding this rank can be reduced to the task of lower bounding the sum $\sum_{i=1}^{q}|w_{S_i}|$, which is established below. 

Choose $w$ from $\{-k,k\}^d$ (where $k = \log n$) uniformly at random. For each $i\in[q]$, let $E_i$ denote the (bad) event that $|w_{S_i}|\leq  \sqrt{r}k/4$. Since $S_i$ is an interval of length $r$, by a standard estimation of binomial coefficients,  
%upon applying the Chernoff bound (\autoref{thm:Chernoff}), 
we obtain that $\Pr[E_i]\leq 1/4$. Then,  by the Chernoff bound (\autoref{thm:Chernoff}), the probability that at least half of the events $E_i$ occur is at most $e^{-\frac{q}{12}}$. Therefore, with probability at least $1 - e^{-\frac{q}{12}}$, 
\[
\sum_{i=1}^q |w_{S_i}| \geq {q\sqrt{r}k/8},
\]
and therefore, by the sub-additivity of $\rk_w(\cdot)$,
\[
\rk_w(A) \leq s^q 2^{-\frac{kq\sqrt{r}}{16}} 
= s^q n^{-\frac{{q\sqrt{r}}}{16}}.
\]

Now, let $\sum_{i=1}^t A_i$ be a $\sbp$ computing $F_{n,d}$ with max-width bounded by $s$, and such that each $A_i$ is ordered set-multilinear with respect to the variable partition $(X_{\sigma_i(1)},\ldots, X_{\sigma_i(d)})$ for some permutation $\sigma_i\in S_d$. %such that $\sum_{i=1}A_i$ computes a polynomial $P(X_1,\ldots, X_d)$ such that $\rk_w(P) = 1$ for every word $w\in\{-k,k\}^d$ with $w_{[d]} = 0$. %Such polynomials exist and we have explicit examples in VNP and VP due to \cite{KushS22} and \cite{KushS23} respectively.
By the union bound and the discussion above, it follows that with probability at least $1-t\cdot e^{-\frac{q}{12}}$, 
\[
\rk_w(F_{n,d}) = \rk_w\left(\sum_{i=1}^t A_i\right) \leq t\cdot s^q n^{-\frac{{q\sqrt{r}}}{16}}.
\]

But now, we can condition on the event that $w_{[d]} = 0$ (which occurs with probability $\Theta(\frac{1}{\sqrt{d}}$)) to establish the existence of a word $w\in \{-k,k\}^d$ with $w_{[d]} = 0$ such that $w$ satisfies $
\rk_w(P)\leq t\cdot s^q n^{-q\sqrt{r}}$. This is because of the given bound $q\geq 12\ln(2t\sqrt{d})$. Because $\rk_w(F_{n,d}) = 1$ for such a $w$ by \autoref{lem:Fn-VP-props}, we conclude that $t\cdot s^q \geq n^{\frac{{q\sqrt{r}}}{16}} \geq n^{\frac{\sqrt{(d-q)q}}{16}}\geq n^{\frac{\sqrt{dq}}{20}}$, where the last inequality follows from our choice of $q$.
%t\geq \left(\frac{n^{\frac{\alpha\sqrt{r}}{4}}}{s}\right)^q.
    
\end{proof}

\section{Tightness of ABP Set-Multilinearization}\label{sec:tightness}

In this section, we prove \autoref{thm:tightness} and in the process, also prove \autoref{thm:large-deg-VBP} and \autoref{thm:small-deg-VBP}. We first establish the following technical lemma that will be essential for these proofs. 

\begin{lem}\label{lem:VBP}
Let $d\leq n$ be growing parameters satisfying $d = \omega(\log n)$. There exist fixed constants $\gamma,c,c_1,c_2>0$ and $\eps\geq 1/100$ such that any $\sbp$ of max-width $s$ and support size $t$ computing $G_{n,d}(X_1,\ldots,X_d,Y_1\ldots,Y_d)$ must satisfy at least one of the following:
\begin{itemize}
    \item either $t\geq 2^{cc_2d}$,
    \item or $t\cdot s^q \geq n^{{{\gamma q(d/q)^\eps}}}$, for any integer $q$ in the range $ [\max\{c_1,(\log t)/c\}, c_2d]$.
\end{itemize}
\end{lem}

\begin{proof}
First, similar to the proof of \autoref{lem:VP}, we observe that for any $\sbp$ computing $G_{n,d}(X,Y)$, we can view each summand as an ordered set-multilinear branching program with respect to only the $X$ variables. In other words, by appropriately collapsing the layers labelled using the $Y$ variables, each summand is a set-multilinear branching program over the field $\F(Y)$ \emph{ordered} with respect to $(X_{\sigma(1)},\ldots, X_{\sigma(d)})$ for some permutation $\sigma\in S_d$. It is easy to see that this collapsing process does not  increase the width or the size of any summand branching program.  The edge labels do get altered however: the coefficients of the $X$ variables in any edge label can now be polynomials in the $Y$ variables (and therefore, field constants in $\F(Y)$).


Let $A$ be such a set-multilinear branching program of width $s$ and depth $d$ that is \emph{ordered} with respect to $(X_1,\ldots, X_d)$\footnote{In general, $A$ may be ordered with respect to an arbitrary permutation $\sigma$, but the assumption that $\sigma$ is the identity permutation in the discussion that follows is without loss of generality.}. Recall that this means that for each $\ell\in[d]$, all edges of the ABP from layer $\ell-1$ to layer $\ell$ are labeled using a linear form {in} $\mathit{X_\ell}$. Given a node $u$ in layer $i$ and a node $v$ in layer $j>i$ of $A$, define $g_{u,v}$ to be the polynomial computed by the ABP restricted to the layers $i+1,\ldots,j-1$ with the source and the sink defined by $u$ and $v$ respectively. Consider an integer $q$ in the range specified in the lemma statement and let $r$ be such that\footnote{\label{foot:rdq} If $q$ does not divide $d$, then we can let $r$ be $\lfloor d/q\rfloor$. In the discussion that immediately follows, we simply bound the relative rank of the `last' component by $1$ and so, the remaining analysis is nearly identical.} $qr = d$. Consider the following decomposition of $A$:
\[
A(X_1,\ldots,X_d) = \sum_{u_1,\ldots, u_{q-1}} \prod_{i = 1}^{q} g_{u_{i-1},u_i},
\]
where $u_0$ and $u_{q}$ are defined to be the source and the sink of $A$ respectively, and for $1\leq i \leq q-1$, $u_i$ varies over all choices of nodes in layer $r\cdot i$. Note that hence, this expression contains at most $s^{q-1}\leq s^q$ terms. Also, note that each $g_{u_{i-1},u_i}$ is set-multilinear over the partition $(X_i)_{i\in S_i}$ where $S_i$ is the set $\{r(i-1) + 1, \ldots, ri\}$\footnote{If $A$ is instead  ordered with respect to $\sigma$, then $S_i$ would be simply defined as the set $\{\sigma(r(i-1) + 1), \ldots, \sigma(ri)\}$.} of length exactly $r$. We now analyze the relative rank of each summand. 

By \autoref{clm:rk-props}, we see that for every appropriate word $w$,
\[
\rk_w\left(\prod_{i = 1}^{q+1} g_{u_{i-1},u_i}\right) = \prod_{i = 1}^{q+1} \rk_{w|_{S_i}}(g_{u_{i-1},u_i})
\leq \prod_{i = 1}^{q} \rk_{w|_{S_i}}(g_{u_{i-1},u_i})\leq\prod_{i = 1}^q  2^{-|w_{S_i}|/2}\leq 2^{-\frac12\left(\sum_{i=1}^q |w_{S_i}|\right)},
\]
from which we observe that the task of upper bounding this rank can be reduced to the task of lower bounding the sum $\sum_{i=1}^{q}|w_{S_i}|$, which is established below. 

Choose $w$ from the distribution $\calD$, as described in Section \ref{subsec:VBP}. We now view the partition $(S_1,\ldots,S_q)$ of $[d]$ as a fixed ``coloring'' of the latter set (and in turn, the $d$-cycle, as described in Section \ref{subsec:VBP}) i.e., each node $i\in[d]$ is assigned the color $k$ if and only if $i\in S_k$. For a pairing $P$ and set $S_k$, define the number of $k$-violations by
$$
V_k(P) = \{P_t \in P : |P_t \cap S_k| = 1\}.
$$
In words, it is the set of pairs in which one color is $k$ and the other color is different from $k$. For some fixed $0<\eps\leq 1/100$, denote
$$
\calG(P) = \{k \in [q] : |V_k(P)| \geq r^{2\eps}\}.
$$

Next, we state a technical lemma that states that with probability exponentially close to $1$, ``many" colors have ``many" violations. The constants $c_1,c_2$ that appear in the statement below indeed define the constants $c_1,c_2$ that are mentioned in the statement of Lemma \ref{lem:VBP}.
\begin{lem}[Many Violations Lemma]\label{lem:many-viol} Let $d\leq n$ be growing parameters satisfying $d = \omega(\log n)$. There exist fixed constants 
$0<\alpha,\beta<1$ and $c_1>0$, $0<c_2\leq 1$ such that for all integers $q$ in the range $ [c_1,c_2d] $ the following holds:
Let $S = (S_1, \ldots, S_q)$ be a partition of the $d$-cycle where each $|S_i| = r$.
Then,
$$
\Pr[|\calG(P)| \leq \alpha q] \leq r^{-\beta q},
$$
where $P\sim \calD_P$.
\end{lem}

We now show, in the claim below, how the preceding lemma can be used to argue that with probability exponentially close to $1$, an arc-partition $w$ exhibits large bias. The constant $c$ that appears below defines the constant $c$ in the statement of Lemma \ref{lem:VBP}.

\begin{claim}\label{clm:high-bias-VBP}
There exists a fixed constant $c>0$ such that
\[   \Pr\left[\sum_{i=1}^q |w_{S_i}| \leq \frac{\alpha qr^\eps k}{64}\right] \leq 2^{-cq},
\]
where the probability is over the choice of $w\sim \calD$.
\end{claim}
\begin{proof}
Let $\calE$ denote the event $\sum_{i=1}^q |w_{S_i}| \leq \frac{\alpha qr^\eps k}{64}$, and $\calA$ denote the event $|\calG(P)| > \alpha q$. From the law of total probability, it follows that $\Pr[\calE]\leq \Pr[\calE|\calA] + \Pr[\overline{\calA}]$, where $\overline{\calA}$ denotes the complement of $\calA$ -- therefore, it suffices to bound $\Pr[\calE|\calA]$. 

Fix a pairing $P\sim \calD_P$ such that the high probability event $\calA$ occurs. Consider an ordering $\sigma$ of the colors in $\calG(P)$. A color $\ell$ is said to be \emph{bright} with respect to an ordering if there are at least $r^{2\eps}/2$ nodes $x$ of color $\ell$ such that the partner of $x$ is colored using a color that appears \emph{after} $\ell$ in the ordering $\sigma$.  Call an ordering $\sigma$ of the nodes in $\calG(P)$ \emph{good} if there are at least $|\calG(P)|/2$ bright colors  with respect to $\sigma$. The observation is that for any ordering $\sigma$ of the colors, either $\sigma$ itself is good, or its reverse is good. We conclude that given any pairing $P$, there exists a good ordering of $\calG(P)$. Fix any such good ordering and let $\calH(P)$ be the collection of bright colors with respect to this ordering. Let the colors in $\calH(P)$ according to this good ordering be $\ell_1,\ldots, \ell_{q'}$.

Next, notice that if the sum $\sum_{j=1}^{q}|w_{S_j}|$ is at most $\frac{\alpha qr^\eps k}{64}$, then so is the sum $\sum_{\ell\in \calH(P)}|w_{S_\ell}|$. Let $q' = |\calH(P)|$ (which is at least $\alpha q/2$ if $|\calG(P)|>\alpha q$). View the sampling of $w$ from $P$ as happening in a specific order, according to the order
of $\ell_1, \ell_2, \ldots, \ell_{q'}$: First define $\Pi$ on pairs with at least one point with color $\ell_1$, then define $\Pi$ on
remaining pairs with at least one point with color $\ell_2$, and so forth. When finished with $\ell_1,\ldots, \ell_{q'}$,
continue to define $\Pi$ on all other pairs.

Observe that if the sum $\sum_{\ell\in \calH(P)}|w_{S_\ell}|$ is at most $\frac{\alpha qr^\eps k}{64}$, then there exists a subset $T\subseteq \calH(P)$ with half its size, $|T|= |\calH(P)|/2$,\footnote{We assume without loss of generality that $|\calH(P)|$ is even to avoid ceilings and floors.} such that every color $\ell$ in $T$ satisfies $|w_{S_\ell}|\leq kr^\eps /16$ (otherwise, we obtain an immediate contradiction). There are at most $2^{|\calH(P)|}$ many choices for such $T$ -- fix such a choice $T$ and relabel the colors in $T$ as $\ell_1,\ldots, \ell_{q''}$, where $q'' = |\calH(P)|/2$ and this order respects the order described in the paragraph above.
For every $\ell_j\in T$, define $E_j$ to be the event that $|w_{S_{\ell_j}}|\leq kr^\eps /16$. 
By choice, conditioned on $E_1,\ldots, E_{j-1}$, there are at least $r^{2\eps}/2$ pairs $P_t$ so that
$|P_t \cap S_{\ell_j}| = 1$ that are not yet assigned a `positive' or `negative' sign (with a magnitude equal to $|k|$). For every such $P_t$, the element in $P_t \cap S_{\ell_j}$
is assigned a positive sign with probability $1/2$, and is independent of any other $P_{t'}$. Therefore, $\Pr[E_j]$ is bounded by the probability that a binomial random variable over a universe of size $r^{2\eps}/2$ lies in any \emph{specific} interval of size $r^\eps/8$. By a standard estimation of binomial coefficients, this probability is bounded by a constant, $1/5$.

Hence, for any fixed choice of $T\subseteq \calH(P)$, for all $\ell_j \in T$,
\[
\Pr[E_j|E_1,\ldots,E_{j-1},P] \leq 1/5.
\]
Therefore, for any fixed choice of $T\subseteq \calH(P)$, by the chain rule, it follows that 
\[
\Pr[\cap_{j\in T}E_j]\leq 5^{-|T|} = 2^{-(\log 5)|\calH(P)|/2}.
\]
There are at most $2^{|\calH(P)|}$ choices for $T$. We conclude that 
\[\Pr[\calE|\calA]\leq 2^{|\calH(P)|} \cdot 2^{-(\log 5)|\calH(P)|/2}\leq 2^{-|\calH(P)|/7}\leq 2^{-\alpha q/14}.
\]
Finally, we note that by the many violations lemma (\autoref{lem:many-viol}), $\Pr[\overline{\calA}]\leq r^{-\beta q}\leq 2^{-\beta q}$. Thus, $\Pr[\calE] \leq \Pr[\calE|\calA] + \Pr[\overline{\calA}] \leq 2^{-\alpha q/14} + 2^{-\beta q} \leq 2^{-cq}$ for some fixed constant $c>0$ which can be defined in terms of $\alpha$ and $\beta$.
\end{proof}

From the claim, it follows that with probability at least $1 - 2^{-cq}$, 
\[
\sum_{i=1}^q |w_{S_i}| \geq \frac{\alpha qr^\eps k}{64},
\]
and therefore, by the sub-additivity of $\rk_w(\cdot)$,
\[
\rk_w(A) \leq s^q 2^{-\frac{\alpha kqr^\eps }{128}} = s^q n^{-{\frac{\alpha qr^\eps }{128}}}.
\]

Now, let $\sum_{i=1}^t A_i$ be a $\sbp$ computing $G_{n,d}$ with max-width bounded by $s$, and such that each $A_i$ is ordered set-multilinear with respect to the variable partition $(X_{\sigma_i(1)},\ldots, X_{\sigma_i(d)})$ for some permutation $\sigma_i\in S_d$. Assume that $2^{cc_1}< t< 2^{cc_2d}$ (if the second inequality does not hold, then the first item in the lemma statement is true; and we deal with the case $t\leq 2^{cc_1}$ at the end of the proof). %such that $\sum_{i=1}A_i$ computes a polynomial $P(X_1,\ldots, X_d)$ such that $\rk_w(P) = 1$ for every word $w\in\{-k,k\}^d$ with $w_{[d]} = 0$. %Such polynomials exist and we have explicit examples in VNP and VP due to \cite{KushS22} and \cite{KushS23} respectively.
By the union bound and the discussion above, it follows that with probability at least $1-t\cdot 2^{-cq}$, 
\[
\rk_w(F_{n,d}) = \rk_w\left(\sum_{i=1}^t A_i\right) \leq t\cdot s^q n^{-{\frac{\alpha qr^\eps }{128}}}.
\]
Since $q> (\log t)/c$ by assumption, note that $1-t\cdot 2^{-cq}>0$. Furthermore, since we are assuming $2^{cc_1}< t$, we have that $q\geq c_1$ and therefore, Lemma \ref{lem:many-viol} applies and so does \autoref{clm:high-bias-VBP} along with the entire discussion above. We conclude that there exists an arc-partition $w\in \{-k,k\}^d$ such that $w$ satisfies $
\rk_w(F_{n,d})\leq t\cdot s^q n^{-{\frac{\alpha qr^\eps }{128}}}$. Because $\rk_w(F_{n,d}) = 1$ for such a $w$ by \autoref{lem:Fn-VBP-props}, we conclude that $t\cdot s^q \geq n^{{\frac{\alpha qr^\eps }{128}}}$.

Finally, if $t$ is bounded by the constant $2^{cc_1}$, then we just add constantly many width-$2$ ordered set-multilinear branching programs which each compute the zero polynomial so that the new support size of the sum becomes larger than $2^{cc_1}$ and the previous case applies, and we are able to conclude $(t + 2^{cc_1})\cdot s^q \geq n^{{\frac{\alpha qr^\eps }{128}}}$. Appropriately defining the constant $\gamma>0$ then lets  us conclude the desired bound $t\cdot s^q \geq n^{{{\gamma q(d/q)^\eps}}}$.
\end{proof}

\begin{proof}[Proof of \autoref{thm:large-deg-VBP}]
Set $\delta = \eps/(1+\eps)$, where $\eps$ is as defined in the Lemma \ref{lem:VBP}. If $t \geq 2^{n^{\delta}}$, then we trivially have that the total-width is at least  ${\exp(\Omega(n^{\delta}))}$, so assume $t \leq 2^{n^{\delta}}$ (so that the second item of \autoref{lem:VBP} applies). We shall show that then, $s = \exp(\Omega(n^{\delta}))$, which will yield the desired result.

By \autoref{lem:VBP}, any $\sbp$ of max-width $s$ and support size $t$ computing $G_{n,n}$ satisfies the inequality $t\cdot s^q \geq n^{{\gamma q(n/q)^\eps}}$, for any integer $q$ in the range $ [(\log t)/c, c_2n]$. Set $q = \lceil n^{\delta}/c\rceil$. Then clearly, $q\geq (\log t)/c$ and so, $q$ lies in the required range.  Plugging in the setting for $q$ in the inequality, we see that 
\[
s \geq \frac{n^{\gamma{({n/q})^\eps}}}{t^{1/q}} \geq \frac{n^{\gamma({cn^{1-\delta}})^{\eps}}}{2^c} = \exp(\Omega(n^{\delta})).
\]
\end{proof}

\begin{proof}[Proof of \autoref{thm:small-deg-VBP}]
We consider cases as follows:
%n^{-{\frac{\alpha q\sqrt[4]{d/q}}{32}}}
\paragraph{\underline{Case $t = \poly(n)$:}} Suppose there is a constant $c'$ such that $t\leq n^{c'}$. Set $q = \lceil(\log n^{c'})/c\rceil = \lceil (c'\log n)/c\rceil$ and note that $q = \Theta(\log n)$. We see that by \autoref{lem:VBP},
\[
s \geq \frac{n^{\gamma{({d/q})^\eps}}}{t^{1/q}}\geq n^{{\gamma (d/q)^\eps}- \frac{c'}{q}}.
\]
Note that $c'/q <1$ and decays to zero as $d$ becomes larger. Furthermore, as $d = \omega(\log n)$ by assumption and $q = \Theta(\log n)$, $d/q = \omega(1)$. We conclude that if $t$ is bounded by a polynomial in $n$, then $s$ must be super-polynomial in $n$.

\paragraph{\underline{Case $s = \poly(n)$:}} Suppose there is a fixed constant $c'\geq 1$ such that $s\leq n^{c'}$. Define $r$ to be the constant $ \max\{(\frac{2c'}{\gamma})^{1/\eps}, \frac{1}{c_2}\}$. Set $q = d/r$ and assume that $t<2^{cq}$. Then $q$ indeed lies in the range to apply \autoref{lem:VBP}: we have $q\leq c_2d$ because $1/r\leq c_2$ by definition, and $q\geq (\log t)/c$ because of the assumption on $t$. We obtain the inequality
\[
t\geq \left(\frac{n^{\gamma ({d/q})^\eps}}{s}\right)^q \geq \left(\frac{n^{{\gamma r^\eps}}}{n^{c'}}\right)^q\geq \left(\frac{n^{\frac{2c'\gamma }{\gamma}}}{n^{c'}}\right)^q = n^{c'q} = n^{\Omega(d)},
\]
which contradicts the assumption that $t<2^{cq}$. Hence, $t \geq 2^{cq} = 2^{\Omega(d)}$ which is indeed super-polynomial in $n$ whenever $d = \omega(\log n)$.

Thus, in either case, it is shown that both $s$ and $t$ cannot be polynomially bounded. Hence, the total-width of any $\sbp$ computing $G_{n,d}$ cannot be polynomially bounded.
\end{proof}

Finally, we observe that \autoref{thm:tightness} follows immediately from (i) the ABP construction of $G_{n,d}$ given in Section \ref{subsec:VBP}, and (ii), the proof of the $s = \poly(n)$ case in \autoref{thm:small-deg-VBP} above.

\appendix
\section{Proof of the Many Violations Lemma}\label{sec:app}
We briefly recall some notation from Section \ref{sec:tightness}. Recall that $w$ is chosen from the distribution $\calD$, as described in Section \ref{subsec:VBP}. For a pairing $P$, and a set $S_k$, we defined the number of $k$-violations as
$$
V_k(P) = \{P_t \in P : |P_t \cap S_k| = 1\}.
$$

In words, it is the set of pairs in which one color is $k$ and the other color is different. We used the following notation to denote the set of colors with ``many'' violations (for some fixed $0<\eps\leq 1/100$)
$$
\calG(P) = \{k \in [q] : |V_k(P)| \geq r^{2\eps}\}.
$$

As mentioned previously, this subsection is adapted from the proof of the (weaker) many violations lemma in \cite{KushS23}, which is in turn adapted from Lemma 4.1 in \cite{DvirMPY12}. 

\paragraph{Lemma.}(Many Violations Lemma restated) \textit{Let $d\leq n$ be growing parameters satisfying $d = \omega(\log n)$. There exist fixed constants 
$0<\alpha,\beta<1$ and $c_1>0$, $0<c_2\leq 1$ such that for all integers $q$ in the range $ [c_1,c_2d] $ the following holds:
Let $S = (S_1, \ldots, S_q)$ be a partition of the $d$-cycle where each $|S_i| = r$.\footnote{As explained in \autoref{foot:rdq} we assume w.l.o.g. $d=rq$.}
Then,
$$
\Pr[|\calG(P)| \leq \alpha q] \leq r^{-\beta q},
$$
where $P\sim \calD_P$.}


\begin{proof}
Set $\alpha = 1/1000$. Fix a partition (or a ``coloring'') $S = (S_1,\ldots, S_q)$ of the $d$-cycle satisfying the conditions of the lemma.
Think of $S$ as a function from the $d$-cycle to the set $[q]$, assigning every node its color in $[q]$. $S(i)$ is the
color of $i$. Use the following definition to partition the proof into cases. For a color $k$, count the
number of jumps in it (with respect to the partition $S$) to be
\[
J_k = \{j \in S_k : k = S(j) \neq S(j + 1)\},
\]
the set of elements $j$ of color $k$ so that $j + 1$ has a color different from $k$. 

\paragraph{Case 1: Many colors with many jumps.}  The high-level idea is that each color with many jumps
has many violations because pairs of the form $(j, j + 1)$ yield violations as soon as they are constructed.

Assume that for at least $q/2$ colors $k$,
$|J_k| > r^{4\eps}$.
Denote by $B \subseteq [q]$ the set of $k$’s that satisfy this inequality. Then, for every $k$ in $B$, there exists
a subset $Q_k \subset J_k$ of size $N = \lceil r^{4\eps}\rceil$. Let
\[
Q \coloneqq \bigcup_{k\in B} Q_k.
\]
We think of the construction of the (random) pairing $P$ as happening in \emph{epochs}, depending on $Q$, as
follows.


For $t > 0$, define the random variable
\[
Q(t) = Q \setminus [L_t - 4, R_t + 4],
\]
the set $Q$ after removing a four-neighborhood of $[L_t, R_t]$. %; if the arc is small enough (i.e., say the distance between $L_t$ and $R_t$ is at most ten), define $Q(t) = Q$. 
For a certain sequence of time steps $t$, we will define special nodes $q_t$ which lie in this small `cloud' around the arc $[L_t,R_t]$ (i.e., within a distance of $4$ on either side of the arc) - it is for these special nodes $q_t$ that the set of pairs $(q_t,q_{t+1})$ will provide many violations. We now formalize this intuition. 

Let $\tau_1 \geq \tau_0 \coloneqq 1$ be the first time $t$ after $\tau_0$ so that the distance between $[L_t
, R_t
]$ and $Q(\tau_0)$ is
at most two. The distance between $[L_{\tau_0
}
, R_{\tau_0}
]$ and $Q(\tau_0)$ is at least five. The size of the arc $[L_t
, R_t
]$ increases by two at each time step. So, $\tau_1 \geq \tau_0+ 2$. Let $q_1$ be an element of $Q(\tau_0)$ that is of distance
at most two from $[L_{\tau_1}
, R_{\tau_1}
]$; if there is more than one such $q_1$, choose arbitrarily. The minimality
of $\tau_1$ implies that $q_1$ is not in $[L_{\tau_1}
, R_{\tau_1}
]$.

Let $\tau_2 \geq \tau_1$ be the first time $t$ after $\tau_1$ so that the distance between $[L_t
, R_t
]$ and $Q(\tau_1)$ is at most
two. Let $q_2$ be an element of $Q(\tau_1)$ that is of distance at most two from $[L_{\tau_2}
, R_{\tau_2}
]$.
Define $\tau_j$, $q_j$ for $j > 2$ similarly, until $Q(\tau_j )$ is empty. As long as $|Q(\tau_j )| \geq 8$, we have $|Q(\tau_j+1)| \geq
|Q(\tau_j )| - 8$. This process, therefore, has at least $qN/16$ steps.
For $1 \leq j \leq qN/16$, denote by $E_j$ the event that during the time between $\tau_j$ and $\tau_{j+1}$ the pair
$\{q_j , q_j + 1\}$ is added to $P$. The pair $\{q_j , q_j + 1\}$ is violating color $S(q_j)$. At time $\tau_j$, even conditioned
on all the past $P_1, \ldots , P_{\tau_j}$, in at most two steps (and before $\tau_{j+1}$) we can add the pair $\{q_j , q_j + 1\}$
to $P$. For every $j$, therefore,
\[
\Pr[E_j |P_1,\ldots , P_{\tau_j}] \geq (1/3)(1/3) = 1/9.
\]
Next, let $N' = \lceil qN/960\rceil$. We want to show that with high probability,  for at least $N'$ many $j$, the event $E_j$ occurs. There are $\binom{\lfloor qN/16\rfloor}{\lceil qN/960 \rceil}$ many ways of choosing a set of indices $j$ of size $N-N'$. Subsequently,
\begin{align*}
\Pr[\text{there is $j_1,\ldots , j_{N'}$ so that $E_{j_1}\cap \cdots \cap E_{j_{N'}}$}] &\geq  1 - \binom{\lfloor qN/16\rfloor}{\lceil qN/960 \rceil}\cdot \left(\frac{8}{9}\right)^{N-N'}\\
&\geq 1 - \left(\frac{960e}{16}\right)^{N'}\cdot \left(\frac{8}{9}\right)^{60N'}\\
&\geq 1 - c^{N'}
\end{align*}
where $0< c < 1$ is a universal constant. Finally, we argue that if there do exist $j_1,\ldots , j_{N'}$ for which the events $E_{j_1},\ldots, E_{j_{N'}}$ occur, then $|\calG(P)|\geq q/1000$. To see this, note that the size of every $Q_k$ is $N$. So, every color $k$ in $B$ can contribute at most $N$ elements to
$j_1, \ldots , j_{N'}$. If $|\calG(P)|<q/1000$, then at most these many colors can contribute more than $r^{2\eps}$ (and up to $N$ elements) - combined, at most $qN/1000$ elements. However, there are at least $q/2 - q/1000$ colors which can contribute only up to $r^{2\eps}$ elements. Again combined, this is not sufficient to cover the $N'$ elements overall (for large enough $d$ and a small enough constant $c_2$ that depends only on $\epsilon$), which is a contradiction. Hence,
\[
\Pr[|\calG(P)| \geq q/1000] \geq \Pr[\text{there is $j_1,\ldots , j_{N'}$ so that $E_{j_1}\cap \cdots \cap E_{j_N'}$}].
\]
and the proof follows in this case as $c^{N'}\ll r^{-\Omega(q)}$.
%\Amir{For this to be the case we must have $r^{4\eps} = N=\Omega(\log d)$. This means that $c_2$ cannot be a constant. Right?}

\paragraph{Case 2: Many colors with few jumps.}
The intuition is that many violations will come
from pairs of the form $\{L_t - 1, R_t + 1\}$ in the construction of the pairing.
Assume that for at least $q/2$ colors $k$,
$|J_k| \leq r^{4\eps}$.
Denote again by $B \subseteq [q]$ the set of $k$’s that satisfy the above inequality. We say that a color $k$ is
noticeable in the arc $A$ if
\[
r^{1 - 6\eps} \leq |S_{k}\cap A| \leq 
|A| - r^{1-6\eps}.
\]

\begin{claim}\label{clm:noticeable}
There are $q' \geq q/2 - 1$ disjoint arcs $A_1,\ldots, A_{q'}$ so that for every $j \in [q'
]$,
\begin{enumerate}
    \item $|A_j | = m = \lfloor r^{1-5\eps}\rfloor$ and,
\item there is a color $k_j$ in $B$ that is noticeable in $A_j$.
\end{enumerate}
Moreover, the colors $k_1,\ldots, k_{q'}$ can be chosen to be pairwise distinct.
\end{claim}

\begin{proof}
For each color $k$ in $B$, there are at least $r$ vertices of color $k$ in the $d$-cycle and at most
$r^{4\eps}$ jumps in the color $k$. Therefore, there is at least one $k$-monochromatic arc of size at least
$r^{1-4\eps}$. Hence, on the $d$-cycle, there are such monochromatic arcs $I_{k_1}
,\ldots, I_{k_{|B|}}$
for the colors
$k_1,\ldots, k_{|B|}$
in $B$, in this order $(1 < 2 <\cdots< D)$.

Consider an arc $A$ of size $m$ included in $I_{k_1}$. Thus $|S_{k_1} \cap A| = m$. If we “slide” the arc $A$ until
it is included in $I_{k_2}$, then $|S_{k_1} \cap A| = 0$. By continuity, there is an intermediate position for the arc
$A$ such that $r^{1-6\eps} \leq |S_{k_1} \cap A| \leq m - r^{1-6\eps}$. This provides the first arc $A_1$ of the claim.

Sliding an arc inside $I_{k_2}$
to inside $I_{k_3}$
shows that there exists an arc $A_2$ such that $r^{1-6\eps}\leq
|S_{k_2} \cap A_2| \leq m - r^{1-6\eps}$. The arcs $A_1$ and $A_2$ are disjoint: The distance of the largest element of $A_1$
and the smallest element of $I_{k_2}$
is at most $m$. The distance of the smallest element of $A_2$ and the
largest element of $I_{k_2}$
is at most $m$. The size of $I_{k_2}$
is larger than $2m$.
Proceed in this way to define $A_3,\ldots, A_{|B|-1}$.
\end{proof}

Use Claim \ref{clm:noticeable} to divide the construction of the (random) pairing into \emph{epochs}. Denote by $A^{(0)}$
the family of arcs given by the claim. Let $\tau_1$ be the first time $t$ that the arc $[L_t
, R_t
]$ hits one of
the arcs in $A^{(0)}$. Denote by $A_1$ that arc that is hit at time $\tau_1$ (break ties arbitrarily). Denote by
$k_1$ the color that is noticeable in $A_1$. Let $\sigma_1$ be the first time $t$ so that $A_1$ is contained in $[L_t
, R_t
]$.
Let $A^{(1)}$ be the subset of $A^{(0)}$ of arcs that have an empty intersection with $[L_{\sigma_1}
, R_{\sigma_1}
]$. Similarly,
let $\tau_2$ be the first time $t$ after $ \sigma_1$ that the arc $[L_t
, R_t
]$ hits one of the arcs in $A^{(1)}$. If there are no
arc in $A^{(1)}$, define $\tau_2 = \infty$. Denote by $A_2$ that arc that is hit at time $\tau_2$. Denote by $k_2$ the color that
is noticeable in $A_2$. 
Let $\sigma_2$ be the first time $t$ so that $A_2$ is contained in $[L_t, R_t]$. 
Let $A^{(2)}$ be the subset of $A^{(1)}$ of arcs that have an empty intersection with $[L_{\sigma_2}, R_{\sigma_2}]$. 
Define $\tau_j , \sigma_j , A_j , k_j, A^{(j)}$ for $j > 2$ analogously.
For every $j \geq 1$, denote by $E_j$ the event that during the time between $\tau_j$ and $\tau_{j+1}$ the number of
pairs added that violate color $k_j$'s at most $r^{2\eps}$. (If $E_j$ does not hold, then $|V_{k_j}
(P)| \geq r^{2\eps}$ and $k_j\in \calG(P)$.) 
The main part of the proof is summarized in the following proposition, whose proof is deferred to
Section \ref{subsec:chessboard}.

\begin{lem}[Chessboard Lemma]\label{lem:chessboard}
There is an absolute constant $0<\eps'\leq 1/100$ such that for every $j \geq 1$, and any choice of pairs $P_1,\ldots,P_{\tau_j}$,
\[\Pr[E_j|P_1,\ldots, P_{\tau_j}, |A^{(j-1)}| \geq 3]\leq r^{-\eps'}.\]
\end{lem}

Given this lemma, let us finish the proof of Lemma \ref{lem:many-viol}. Define $q'' = \lfloor q'/10 \rfloor$ and let $T$ denote the event that the number of $j$'s for which $|A^{(j)}|\geq 3$ is at least $q''$. First, we argue that $T$ occurs with high probability. 

For any $j\geq 1$, consider the evolution of the arc $[L_t,R_t]$ between the time steps $\tau_j$ (when it first hits arc $A_j$) and $\sigma_j$ (when it completely engulfs it). During this epoch, let us call the evolution of $[L_t,R_t]$ \emph{in} the `direction' of $A_j$ as \emph{good} (labelled `$G$') and \emph{away} from the direction of $A_j$ as \emph{bad} (`$B$'). To this end, for any time step in this epoch, we can \emph{code} the three possible choices for the evolution of $[L_t,R_t]$ as $GG$ (when the arc is grown \emph{in} the direction of $A_j$), $GB$ (when it is grown equally on either side), or $BB$ (when it is grown \emph{away} from the direction of $A_j$). Consequently, the evolution of $[L_t,R_t]$ during this epoch can be realized as a sequence consisting of the symbols $G$ and $B$.

Consider the sequence $s$ of $G$'s and $B$'s obtained by concatenating the sequences corresponding to all the epochs (ignoring the choices made at time steps that do \emph{not} lie in such epochs, i.e., between $\tau_j$ and $\sigma_j$ for some $j$ - as there is no corresponding notion of a `good' direction outside such epochs). The intuition is that if $|A^{(q'')}| < 3$ (i.e., if $T$ does not occur), then there must be an extremely large number of $B$'s compared to $G$'s (i.e., the arc $[L_t,R_t]$ evolves disproportionately in the \emph{bad} direction) in the concatenated string $s$, which should occur only with a vanishingly small probability.

Consider the sub-string $s'$ of $s$ that corresponds to the choices made only for the nodes in $A^{(0)}\setminus A^{(q'')}$. Note that there are precisely $mq''$ many $G$'s in $s'$. Suppose $|A^{(q'')}| = 2$ for concreteness (the cases $|A^{(q'')}| = 1$ and $|A^{(q'')}| = 0$ are similar). This implies that there are $m(q' - 2 - q'')$ many $B$'s in $s'$. Since only up to $mq''$ many of these $B$'s may appear as a result of the evolution making a choice of the form $GB$, it follows that the evolution of $[L_t,R_t]$ must make a choice of the form $BB$ at least $m(q' - 2 - 2q'')/2$ times out of a possible $m(q' - 2)/2$, in order to cover the elements of $A^{(0)}\setminus A^{(q'')}$. Denote $q_1 \coloneqq (q' - 2)/2$. By the union bound, this probability is at most 
\[
\Pr[|A^{(q'')}| = 2] \leq \binom{mq_1}{mq''}\cdot \left(\frac{1}{3}\right)^{m(q_1 - q'')} < c_2^{mq''}
\]
for some universal constant $0<c_2<1$. Similarly, we have bounds for both $\Pr[|A^{(q'')}| = 1]$ and $\Pr[|A^{(q'')}| = 0]$ and it follows that $\Pr[T] \geq 1 - c^{mq''}$ for some universal constant $0<c<1$. 

\begin{remark}
The argument above for showing that $T$ occurs with high probability differs considerably from \cite{DvirMPY12}, where the corresponding event is sketched to occur with probability only at least $1-dc^{m^{1/3}}$, which is not strong enough for our purposes.  
\end{remark}

Next, note that

\[
 \Pr[|\calG(P)| < q/1000] \leq \Pr[|\calG(P)|< q/1000 \cap T] + \Pr[\neg T]\leq \Pr[|\calG(P)|< q/1000 | T] + \Pr[\neg T].
\]
If $|\calG(P)| < q/1000$, then at least $q/2 - q/1000$ colors in $B$ have at most $r^{2\eps}$ many violations. Since $q'' = \lfloor q'/10\rfloor < q/2 - q/1000$, in particular, there must exist at least $q''/2$ colors within the \emph{first} $q''$ colors (here we are using the ordering of colors as provided by Claim \ref{clm:noticeable}) for which there are at most $r^{2\eps}$ many violations. We then obtain the following by conditioning on $T$, using the union bound.
\[
 \Pr[|\calG(P)|< q/1000 | T] \leq 2^{q''} \max_{H = \{j_1<\cdots< j_{q''/2}\} \subset [q'']} \Pr[E_{j_1},\ldots, E_{j_{q''/2}}||A^{(q'')}|\geq 3]
\]
For a fixed choice of $H$, by the chain rule and Lemma \ref{lem:chessboard}, we have
\begin{align*}
\Pr[E_{j_1}\cap \cdots \cap E_{j_{q''/2}}||A^{(q'')}|\geq 3] &= \Pr[E_{j_1}|T]\cdot \Pr[E_{j_2}|E_{j_1}\cap T]\cdot \cdots\cdot  \Pr[E_{j_{q''/2}}|E_{j_{q''/2 - 1}}\cap \cdots \cap E_{j_1} \cap T]\\ &\leq r^{-\eps' q''/2}\leq r^{-\eps' q'/20} \leq r^{-q/400}.
\end{align*}
%\Amir{It seems to me that using \autoref{lem:chessboard} the bound above should be only $\gamma^{O(q)}$. }

Overall, setting $\beta$ appropriately, we conclude that 
\[
\Pr[|\calG(P)|< \alpha q] \leq r^{-\beta q}.
\]
\end{proof}

\subsection{Proof of the Chessboard Lemma}\label{subsec:chessboard}

To prove Lemma \ref{lem:chessboard}, we use a different point of view of the random process. We begin by describing this different view, and later describe its formal connection to the distribution on pairings. This subsection is adapted from Section 5 of \cite{DvirMPY12} and closely follows their argument, though with numerous parameter changes to suit our demands.

The view uses two definitions. One is a standard definition of a two-dimensional random
walk, and the other is a definition of a “chessboard” configuration in the plane. The proof of the
proposition will follow by analyzing the behavior of the random walk on the “chessboard”.
Let $d$ be as above and $m$ be as defined in Lemma \ref{clm:noticeable}. The random walk $W$ on $\N^2$
is defined as follows. It
starts at the origin, $W_0 = (0, 0)$. At every step it move to one of three nodes, independently of
previous choices,

\[
W_{t+1} = \begin{cases}
     W_t + (0,2) &\text{with probability $1/3$}\\
     W_t + (1,1) &\text{with probability $1/3$}\\
     W_t + (2,0) &\text{with probability $1/3$}
\end{cases}
\]

At time $t$, the $L_1$-distance of $W_t$ from the origin is thus $2t$.

The “chessboard” is defined as follows. Let $\alpha_1 : [m] \rightarrow \{0, 1\}$ and $\alpha_2 : [2m] \rightarrow \{0, 1\}$ be two
Boolean functions. The functions $\alpha_1, \alpha_2$ induce a “chessboard” structure on the board $[m] \times [2m]$.
A position in the board $\xi = (\xi_1, \xi_2)$ is colored either white or black. It is colored black if $\alpha_1(\xi_1) \neq
\alpha_2(\xi_2)$ and white if $\alpha_1(\xi_1) = \alpha_2(\xi_2)$. We say that the “chessboard” is well-behaved if

\begin{enumerate}
    \item $\alpha_1$ is far from constant:
    \[
r^{1-6\eps} \leq
|\{\xi_1 \in [m] : \alpha_1(\xi_1) = 1\}|  \leq m - r^{1-6\eps}.
    \]
    \item $\alpha_1$ does not contain many jumps:
    \[
|\{\xi_1 \in [m - 1] : \alpha_1(\xi_1) \neq \alpha_1(\xi_1 + 1)\}|\leq r^{4\eps}
    \]
    \item $\alpha_2$ does not contain many jumps:
   \[
|\{\xi_2 \in [2m - 1] : \alpha_2(\xi_2) \neq \alpha_2(\xi_2 + 1)\}|\leq r^{4\eps}
    \]
\end{enumerate}

Consider a random walk $W$ on top of the “chessboard” and stop it when reaching the boundary
of the board (i.e., when it tries to make a step outside the board $[m] \times [2m]$). We define a good
step to be a step of the form $(1, 1)$ that lands in a black block. We will later relate good steps to
violating edges. Our goal is, therefore, to show that a typical $W$ makes many good steps.

\begin{lem}\label{lem:good-steps}
    Assume the chessboard is well-behaved. There is a constant $0<\eps'\leq 1/100$ such that the probability that $W$ makes less than $r^{4\eps}$
good steps is at most $r^{-\eps'}$.
\end{lem}

We use this lemma to show Lemma \ref{lem:chessboard}.

\begin{proof}[Proof of Lemma \ref{lem:chessboard} given Lemma \ref{lem:good-steps}]
    Recall that $A_j$ is an arc of size $|A_j| = m = \lfloor r^{1-\delta - \delta'}\rfloor$
 so that there is a color
$k_j$ satisfying
\begin{equation}\label{eq:noticeable}
r^{1-6\eps} \leq |S_{k}\cap A| \leq 
|A| - r^{1-6\eps}.
\end{equation}

Furthermore, condition on $P_1, \ldots, P_{\tau_j}$
, $|A^{(j-1)}| \geq 3$. Assume without loss of generality that $R_{\tau_j}$
is in $A_j$ (when $L_{\tau_j}$
is
in $A_j$, the analysis is similar). The distance of $R_{\tau_j}$
from the smallest element of $A_j$ is at most one
(the length of “one step to the right” is between zero and two). We now grow the random interval
until $\sigma_j$, i.e., as long as $R_t$ stays in $A_j$. At the same time, $L_t$ performs a movement to the left.
Since $|A^{(j-1)}| \geq 3$, there are at least $2m$ steps for $L_t$ to take to the left before hitting $A_j$.
There is a one-to-one correspondence between pairings $P$ and random walks $W$ using the correspondence
\[
P_{t+1} = \{L_t - 2, L_t - 1\} \longleftrightarrow W_{t+1} = W_t + (0, 2),
\]
\[
P_{t+1} = \{L_t - 1, R_t + 1\} \longleftrightarrow W_{t+1} = W_t + (1, 1), \]
\[
P_{t+1} = \{R_t + 1, R_t + 2\} \longleftrightarrow W_{t+1} = W_t + (2, 0).
\]

Define the function $\alpha_1$ to be $1$ at positions of $A_j$ with color $k_j$, and $0$ at the other positions. Set
the function $\alpha_2$ as to describe the color $k_j$ from $L_{\tau_j}$
leftward. The “chessboard” is well-behaved by
(\ref{eq:noticeable}) and since $k_j$ is in the set $B$ defined in case 2 of the proof of Lemma \ref{lem:many-viol} (so there are not many
jumps for the color $k_j$).
Finally, if $W$ makes a good step, then the corresponding pair added to $P$ violated color $k_j$ .
So, if $E_j$ holds for $P$, then the corresponding $W$ makes less than $r^{4\eps}$ good steps. Formally, by
Lemma \ref{lem:many-viol},

\[
\Pr[E_j|P_1,\ldots, P_{\tau_j}, |A^{(j-1)}| \geq 3]\leq \Pr[\text{$W$ makes less than $r^{4\eps}$ good steps}] \leq r^{-\eps'}.
\]
\end{proof}

\begin{proof}[Proof of Lemma \ref{lem:good-steps}]
    Define three events $E_R, E_C, E_D$, all of which happen with small probability, so that every $W$ that
is not in their union makes many good steps.

Call a subset of the board of the form $I \times [2m]$ or $[m] \times I$, where $I$ is a sub-interval, a \emph{region}.
The width of a region is the size of $I$. Let $R$ be the set of regions of width at least $r^{8\eps}$. The size
of $R$ is at most $2m^2$. For a region $t$ in $R$, denote by $E_t$ the event that the number of steps of the
form $(1, 1)$ that $W$ makes in $t$ is less than $r^{4\eps}$ \emph{given} that it makes at least $r^{6\eps}$ steps in $r$. Denote

\[
E_R = \bigcup_{t\in R} E_t
\]

To estimate the probability of $E_t$, note that we can simply apply the Chernoff bound to a sum of $r^{6\eps}$ Bernoulli random variables with $p = 1/3$. By the union bound, we conclude that there is a universal constant $0<c<1$ such that

\[
\Pr[E_R] \leq c^{r^{6\eps}}.
\]

Denote by $H$ the set of all points in the board with $L_1$-norm at least $m^{5/8}$. At time $T$ the
random walk $W$ is distributed along all points in $\N^2$ of $L_1$-norm exactly $T$. The distribution of
$W$ on this set is the same as that of a random walk on $\Z$ that is started at $0$, and moves at every
step to the right with probability $1/3$, stays in place with probability $1/3$ and moves to the left with probability $1/3$. The probability
that such a random walk on $\Z$ is at a specific point in $\Z$ at time $T$ is at most $O(T^{-1/2})$. Hence, for
every point $h$ in $H$,

\[
\Pr[W \text{ hits } h]\leq O(m^{-5/16}).
\]

Call a point $c = (\xi_1, \xi_2)$ in the board a corner if both $(\xi_1, \xi_2)$ and $(\xi_1 + 1, \xi_2 + 1)$ are of the same
color $\kappa \in \{\mathrm{black}, \mathrm{white}\}$, but $(\xi_1 + 1, \xi_2)$ and $(\xi_1, \xi_2 + 1)$ are not of color $\kappa$. For a corner $c$, denote
by $\Delta(c)$ the $r^{8\eps}$-neighborhood of $c$ in $L_1$-metric. Denote by $\Delta$ the union over all $\Delta(c)$, for corners
$c$ in $H$. Denote by $E_C$ the event that $W$ hits any point in $\Delta$. Since the board is well-behaved, the
number of jumps in each of $\alpha_1, \alpha_2$ is at most $r^{4\eps}$. Therefore, the number of corners is at most
$r^{8\eps}$. By the union bound,
\[
\Pr[E_C] \leq O(r^{8\eps}r^{16\eps}m^{-5/16}) \leq r^{-\eps'},
\]
where in the last step, we used $m \geq r^{1 - 5\eps}$. Note that plugging in, say, $\eps = \eps' = 1/100$ indeed makes the inequality true. Next, let $m' = \lceil m^{5/8}\rceil$. Define three (vertical) lines: $D_1$ is the line $\{m'\} \times [2m]$, $D_2$ is the line
$\{2m'\} \times [2m]$ and $D_3$ is the line $\{m - m'\} \times [2m]$. Denote by $E_D$ the event that $W$ does not cross
the line $D_3$ before stopped (i.e., hitting the boundary of the board). Chernoff’s bound implies that there is a universal constant $0<c<1$ for which
\[
P[E_D] \leq c^m.
\]

To conclude the proof by the union bound, it suffices to show that for every $W$ not in $E_R \cup
E_C \cup E_D$, the walk $W$ makes at least $r^{4\eps}$ good steps. Fix such a walk $W$. Since $W \notin E_D$, we
know that $W$ crosses the line $D_2$. 


We consider several cases. Define a \emph{block} to be a maximal monochromatic rectangle in the board.
The board is thus partitioned into black blocks and white blocks - which is what led \cite{DvirMPY12} to calling it a  “chessboard.”
We now think of the board $[m] \times [2m]$ as drawn in the plane with $(1, 1)$ at the bottom-left corner and
$(m, 2m)$ at the upper-right corner.

\textbf{Case 1:}  The walk $W$ does not hit any white block after crossing $D_1$ and before crossing $D_2$.
In this case, all steps taken in the \emph{region} whose left border is $D_1$ and right border is $D_2$ are in a black
area. The number of such steps is at least $m^{5/8}/2\gg r^{6\eps}$. Since $W \notin E_R$, the claim holds.

\textbf{Case 2:}  The walk W hits a white block after crossing $D_1$ and before crossing $D_2$. Let us label the blocks as follows: we associate every block with a pair $\langle \eta_1, \eta_2\rangle$ where $\eta_1$ is between $1$ and the number of jumps in $\alpha_1$ and $\eta_2$ is between $1$ and the number of jumps
in $\alpha_2$. So, the label of the ``bottom-left'' is $\langle 1,1\rangle$, the label of the block ``above'' it is $\langle 1,2\rangle$ and the label of the block ``to its right'' is $\langle 2,1\rangle$, etc.  There are
two sub-cases to consider: 

\textbf{Sub-case 1:} At some point after crossing $D_1$ and before crossing $D_3$, there are two white blocks
of the form $\langle \eta_1, \eta_2\rangle$, $\langle\eta_1 + 1, \eta_2 + 1\rangle$ so that $W$ intersects
both blocks. Let $c$ be the corner between these two blocks (which must exist by definition). Since
$W \notin E_C$, we know that $W$ does not visit $\Delta(c)$. Therefore, $W$ must walk in a black area around
$\Delta(c)$. Every path surrounding $\Delta(c)$ has length at least $r^{8\eps}$. Since $W \notin E_R$, the claim holds.

\textbf{Sub-case 2:} At all times after crossing $D_1$ and before crossing $D_3$, the walk never moves from a
white block $\langle \eta_1, \eta_2 \rangle$ to one of the two white blocks $\langle\eta_1 + 1, \eta_2 + 1\rangle$, $\langle \eta_1 - 1, \eta_2 - 1\rangle$. Since $W \notin E_D$,
this is indeed the last case. The width of a combinatorial rectangle in the board is the size of its
“bottom side” (i.e., the corresponding subset of $[m]$). Let $\eta$ be the first white block $W$ hits after
crossing $D_1$. Let $\Sigma$ be the family of black blocks that are to the right but on the same height as $\eta$.
Define $Z$ as the maximal width of a rectangle of the form $\sigma \cap [0, m - m_0 - 1] \times [2m]$ over all $\sigma \in \Sigma$.
Since the board is well-behaved, it follows (from the first condition) that the
total width of the black area on the same height as $\eta$ is at least $r^{1-6\eps}$. Also, since we are in case 2, the left border of $\eta$ is to the left of $D_2$. Therefore, the
total width of the black area to the right of the left border of $\eta$ and to the
left of $D_3$, on the same height as $\eta$ is at least $r^{1-6\eps}- 3m'$. Therefore, since the number of jumps is at most $r^{4\eps}$,
\[
Z\geq (r^{1-6\eps}- 3m')/r^{4\eps} \gg r^{8\eps}.
\]

Since we are in this sub-case, the walk $W$ must “go through” every black block it hits: it can go from
bottom side to upper side or from left side to right side (but not from left side to upper side or
from bottom side to right side). Consider the behaviour of $W$ after it hits $\eta$: starting from a white block, because $W \notin E_D$, it is guaranteed to cross $D_3$. Therefore, the color of the block that $W$ ``exits'' from from each \emph{column} must keep alternating between white and black. For each black block in $\Sigma$, therefore, there exists
a black block in the same column that $W$ crosses horizontally. Focusing on one such black block
of width $Z$, since $W \notin E_R$, the claim holds.
\end{proof}

\newcommand{\etalchar}[1]{$^{#1}$}
\begin{thebibliography}{DMPY12}

\bibitem[AFS{\etalchar{+}}18]{AndersonFSSV18}
Matthew Anderson, Michael~A. Forbes, Ramprasad Saptharishi, Amir Shpilka, and
  Ben~Lee Volk.
\newblock Identity testing and lower bounds for read-\emph{k} oblivious
  algebraic branching programs.
\newblock {\em {ACM} Trans. Comput. Theory}, 10(1):3:1--3:30, 2018.

\bibitem[AGKS15]{AgrawalGKS15}
Manindra Agrawal, Rohit Gurjar, Arpita Korwar, and Nitin Saxena.
\newblock Hitting-sets for {ROABP} and sum of set-multilinear circuits.
\newblock {\em {SIAM} J. Comput.}, 44(3):669--697, 2015.

\bibitem[AR16]{ArvindR16}
Vikraman Arvind and S.~Raja.
\newblock Some lower bound results for set-multilinear arithmetic computations.
\newblock {\em Chic. J. Theor. Comput. Sci.}, 2016, 2016.

\bibitem[BDS22]{mfcs/BhargavDS22}
C.~S. Bhargav, Sagnik Dutta, and Nitin Saxena.
\newblock Improved lower bound, and proof barrier, for constant depth algebraic
  circuits.
\newblock In Stefan Szeider, Robert Ganian, and Alexandra Silva, editors, {\em
  47th International Symposium on Mathematical Foundations of Computer Science,
  {MFCS} 2022, August 22-26, 2022, Vienna, Austria}, volume 241 of {\em
  LIPIcs}, pages 18:1--18:16. Schloss Dagstuhl - Leibniz-Zentrum f{\"{u}}r
  Informatik, 2022.

\bibitem[BDS24]{BDS23}
C.S. Bhargav, Prateek Dwivedi, and Nitin Saxena.
\newblock Lower bounds for the sum of small-size algebraic branching programs.
\newblock {\em To appear in the proceedings of the Annual Conference on Theory
  and Applications of Models of Computation}, 2024.
\newblock \url{https://www.cse.iitk.ac.in/users/nitin/papers/sumRO.pdf}.

\bibitem[BG22]{BhargavaG22}
Vishwas Bhargava and Sumanta Ghosh.
\newblock Improved hitting set for orbit of roabps.
\newblock {\em Comput. Complex.}, 31(2):15, 2022.

\bibitem[BS21]{BishtS21}
Pranav Bisht and Nitin Saxena.
\newblock Blackbox identity testing for sum of special roabps and its border
  class.
\newblock {\em Comput. Complex.}, 30(1):8, 2021.

\bibitem[B{\"{u}}r00]{burg}
Peter B{\"{u}}rgisser.
\newblock Cook's versus valiant's hypothesis.
\newblock {\em Theor. Comput. Sci.}, 235(1):71--88, 2000.

\bibitem[CKSV22]{ChatterjeeKSV22}
Prerona Chatterjee, Mrinal Kumar, Adrian She, and Ben~Lee Volk.
\newblock Quadratic lower bounds for algebraic branching programs and formulas.
\newblock {\em Comput. Complex.}, 31(2):8, 2022.

\bibitem[DMPY12]{DvirMPY12}
Zeev Dvir, Guillaume Malod, Sylvain Perifel, and Amir Yehudayoff.
\newblock Separating multilinear branching programs and formulas.
\newblock In Howard~J. Karloff and Toniann Pitassi, editors, {\em Proceedings
  of the 44th Symposium on Theory of Computing Conference, {STOC} 2012, New
  York, NY, USA, May 19 - 22, 2012}, pages 615--624. {ACM}, 2012.

\bibitem[FS13]{ForbesS13}
Michael~A. Forbes and Amir Shpilka.
\newblock Quasipolynomial-time identity testing of non-commutative and
  read-once oblivious algebraic branching programs.
\newblock In {\em 54th Annual {IEEE} Symposium on Foundations of Computer
  Science, {FOCS} 2013, 26-29 October, 2013, Berkeley, CA, {USA}}, pages
  243--252. {IEEE} Computer Society, 2013.

\bibitem[FSS14]{ForbesSS14}
Michael~A. Forbes, Ramprasad Saptharishi, and Amir Shpilka.
\newblock Hitting sets for multilinear read-once algebraic branching programs,
  in any order.
\newblock In David~B. Shmoys, editor, {\em Symposium on Theory of Computing,
  {STOC} 2014, New York, NY, USA, May 31 - June 03, 2014}, pages 867--875.
  {ACM}, 2014.

\bibitem[GG20]{guo2020improved}
Zeyu Guo and Rohit Gurjar.
\newblock Improved explicit hitting-sets for roabps.
\newblock In {\em Approximation, Randomization, and Combinatorial Optimization.
  Algorithms and Techniques (APPROX/RANDOM 2020)}. Schloss
  Dagstuhl-Leibniz-Zentrum f{\"u}r Informatik, 2020.

\bibitem[GKS17]{GurjarKS17}
Rohit Gurjar, Arpita Korwar, and Nitin Saxena.
\newblock Identity testing for constant-width, and any-order, read-once
  oblivious arithmetic branching programs.
\newblock {\em Theory Comput.}, 13(1):1--21, 2017.

\bibitem[GKST17]{GurjarKST17}
Rohit Gurjar, Arpita Korwar, Nitin Saxena, and Thomas Thierauf.
\newblock Deterministic identity testing for sum of read-once oblivious
  arithmetic branching programs.
\newblock {\em Comput. Complex.}, 26(4):835--880, 2017.

\bibitem[GR21]{GR21}
Purnata Ghosal and B.~V.~Raghavendra Rao.
\newblock Limitations of sums of bounded read formulas and abps.
\newblock In Rahul Santhanam and Daniil Musatov, editors, {\em Computer Science
  - Theory and Applications - 16th International Computer Science Symposium in
  Russia, {CSR} 2021, Sochi, Russia, June 28 - July 2, 2021, Proceedings},
  volume 12730 of {\em Lecture Notes in Computer Science}, pages 147--169.
  Springer, 2021.

\bibitem[GV20]{GurjarV20}
Rohit Gurjar and Ben~Lee Volk.
\newblock Pseudorandom bits for oblivious branching programs.
\newblock {\em {ACM} Trans. Comput. Theory}, 12(2):8:1--8:12, 2020.

\bibitem[Jan08]{Jansen08}
Maurice~J. Jansen.
\newblock Lower bounds for syntactically multilinear algebraic branching
  programs.
\newblock In Edward Ochmanski and Jerzy Tyszkiewicz, editors, {\em Mathematical
  Foundations of Computer Science 2008, 33rd International Symposium, {MFCS}
  2008, Torun, Poland, August 25-29, 2008, Proceedings}, volume 5162 of {\em
  Lecture Notes in Computer Science}, pages 407--418. Springer, 2008.

\bibitem[Kay12]{Kayal12}
Neeraj Kayal.
\newblock An exponential lower bound for the sum of powers of bounded degree
  polynomials.
\newblock {\em Electron. Colloquium Comput. Complex.}, {TR12-081}, 2012.

\bibitem[KLSS17]{KayalLSS17}
Neeraj Kayal, Nutan Limaye, Chandan Saha, and Srikanth Srinivasan.
\newblock An exponential lower bound for homogeneous depth four arithmetic
  formulas.
\newblock {\em {SIAM} J. Comput.}, 46(1):307--335, 2017.

\bibitem[KNS20]{KayalNS20}
Neeraj Kayal, Vineet Nair, and Chandan Saha.
\newblock Separation between read-once oblivious algebraic branching programs
  (roabps) and multilinear depth-three circuits.
\newblock {\em {ACM} Trans. Comput. Theory}, 12(1):2:1--2:27, 2020.

\bibitem[KS17]{/KumarS17}
Mrinal Kumar and Shubhangi Saraf.
\newblock On the power of homogeneous depth 4 arithmetic circuits.
\newblock {\em {SIAM} J. Comput.}, 46(1):336--387, 2017.

\bibitem[KS22]{KushS22}
Deepanshu Kush and Shubhangi Saraf.
\newblock Improved low-depth set-multilinear circuit lower bounds.
\newblock In Shachar Lovett, editor, {\em 37th Computational Complexity
  Conference, {CCC} 2022, July 20-23, 2022, Philadelphia, PA, {USA}}, volume
  234 of {\em LIPIcs}, pages 38:1--38:16. Schloss Dagstuhl - Leibniz-Zentrum
  f{\"{u}}r Informatik, 2022.

\bibitem[KS23]{KushS23}
Deepanshu Kush and Shubhangi Saraf.
\newblock Near-optimal set-multilinear formula lower bounds.
\newblock In Amnon Ta{-}Shma, editor, {\em 38th Computational Complexity
  Conference, {CCC} 2023, July 17-20, 2023, Warwick, {UK}}, volume 264 of {\em
  LIPIcs}, pages 15:1--15:33. Schloss Dagstuhl - Leibniz-Zentrum f{\"{u}}r
  Informatik, 2023.

\bibitem[KSS14]{KayalSS14}
Neeraj Kayal, Chandan Saha, and Ramprasad Saptharishi.
\newblock A super-polynomial lower bound for regular arithmetic formulas.
\newblock In David~B. Shmoys, editor, {\em Symposium on Theory of Computing,
  {STOC} 2014, New York, NY, USA, May 31 - June 03, 2014}, pages 146--153.
  {ACM}, 2014.

\bibitem[KST16]{KayalST16}
Neeraj Kayal, Chandan Saha, and S{\'{e}}bastien Tavenas.
\newblock An almost cubic lower bound for depth three arithmetic circuits.
\newblock In Ioannis Chatzigiannakis, Michael Mitzenmacher, Yuval Rabani, and
  Davide Sangiorgi, editors, {\em 43rd International Colloquium on Automata,
  Languages, and Programming, {ICALP} 2016, July 11-15, 2016, Rome, Italy},
  volume~55 of {\em LIPIcs}, pages 33:1--33:15. Schloss Dagstuhl -
  Leibniz-Zentrum f{\"{u}}r Informatik, 2016.

\bibitem[LST21]{LST1}
Nutan Limaye, Srikanth Srinivasan, and S{\'{e}}bastien Tavenas.
\newblock Superpolynomial lower bounds against low-depth algebraic circuits.
\newblock In {\em 62nd {IEEE} Annual Symposium on Foundations of Computer
  Science, {FOCS} 2021, Denver, CO, USA, February 7-10, 2022}, pages 804--814.
  {IEEE}, 2021.

\bibitem[LST22]{LST3}
Nutan Limaye, Srikanth Srinivasan, and S{\'{e}}bastien Tavenas.
\newblock On the partial derivative method applied to lopsided set-multilinear
  polynomials.
\newblock In Shachar Lovett, editor, {\em 37th Computational Complexity
  Conference, {CCC} 2022, July 20-23, 2022, Philadelphia, PA, {USA}}, volume
  234 of {\em LIPIcs}, pages 32:1--32:23. Schloss Dagstuhl - Leibniz-Zentrum
  f{\"{u}}r Informatik, 2022.

\bibitem[MU05]{mitzenmacher_upfal_2005}
Michael Mitzenmacher and Eli Upfal.
\newblock {\em Probability and Computing: Randomized Algorithms and
  Probabilistic Analysis}.
\newblock Cambridge University Press, 2005.

\bibitem[Nis91]{Nisan91}
Noam Nisan.
\newblock Lower bounds for non-commutative computation (extended abstract).
\newblock In Cris Koutsougeras and Jeffrey~Scott Vitter, editors, {\em
  Proceedings of the 23rd Annual {ACM} Symposium on Theory of Computing, May
  5-8, 1991, New Orleans, Louisiana, {USA}}, pages 410--418. {ACM}, 1991.

\bibitem[NW97]{NisanW97}
Noam Nisan and Avi Wigderson.
\newblock Lower bounds on arithmetic circuits via partial derivatives.
\newblock {\em Comput. Complex.}, 6(3):217--234, 1997.

\bibitem[Raz06]{Raz06}
Ran Raz.
\newblock Separation of multilinear circuit and formula size.
\newblock {\em Theory Comput.}, 2(6):121--135, 2006.

\bibitem[Raz13]{Raz-Tensor}
Ran Raz.
\newblock Tensor-rank and lower bounds for arithmetic formulas.
\newblock {\em J. {ACM}}, 60(6):40:1--40:15, 2013.

\bibitem[RR20]{RR20}
C.~Ramya and B.~V.~Raghavendra Rao.
\newblock Lower bounds for special cases of syntactic multilinear abps.
\newblock {\em Theor. Comput. Sci.}, 809:1--20, 2020.

\bibitem[RS05]{RazS05}
Ran Raz and Amir Shpilka.
\newblock Deterministic polynomial identity testing in non-commutative models.
\newblock {\em Comput. Complex.}, 14(1):1--19, 2005.

\bibitem[RY08]{RazY08}
Ran Raz and Amir Yehudayoff.
\newblock Balancing syntactically multilinear arithmetic circuits.
\newblock {\em Comput. Complex.}, 17(4):515--535, 2008.

\bibitem[RY09]{RazY09}
Ran Raz and Amir Yehudayoff.
\newblock Lower bounds and separations for constant depth multilinear circuits.
\newblock {\em Comput. Complex.}, 18(2):171--207, 2009.

\bibitem[Sap15]{Saptarishi-survey}
Ramprasad Saptharishi.
\newblock A survey of lower bounds in arithmetic circuit complexity.
\newblock {\em Github Survey}, 2015.

\bibitem[SY10]{ShpilkaY10}
Amir Shpilka and Amir Yehudayoff.
\newblock Arithmetic circuits: {A} survey of recent results and open questions.
\newblock {\em Found. Trends Theor. Comput. Sci.}, 5(3-4):207--388, 2010.

\bibitem[TLS22]{LST2}
S{\'{e}}bastien Tavenas, Nutan Limaye, and Srikanth Srinivasan.
\newblock Set-multilinear and non-commutative formula lower bounds for iterated
  matrix multiplication.
\newblock In Stefano Leonardi and Anupam Gupta, editors, {\em {STOC} '22: 54th
  Annual {ACM} {SIGACT} Symposium on Theory of Computing, Rome, Italy, June 20
  - 24, 2022}, pages 416--425. {ACM}, 2022.

\bibitem[Val79]{Val79}
L.~G. Valiant.
\newblock Completeness classes in algebra.
\newblock In {\em Proceedings of the Eleventh Annual ACM Symposium on Theory of
  Computing}, STOC '79, page 249–261, New York, NY, USA, 1979. Association
  for Computing Machinery.

\end{thebibliography}
\end{document}
